% This chapter is a faithful blueprint adaptation of the indicated part of
% Cao--Zhu's revised TeX source. Source-line sentinels preserve auditability.

\chapter{Long Time Behaviors}
\label{chap:cz-long-time-behavior}

\bigskip
Let $M$ be a complete manifold of dimension $n$. Consider a
solution of the Ricci flow $g_{ij}(x,t)$ on $M$ and on a maximal
time interval $[0,T)$. When $M$ is compact, we usually consider
the \textbf{normalized Ricci flow} \index{normalized Ricci flow}
$$
\frac{\partial g_{ij}}{\partial t}=\frac{2}{n}rg_{ij}-2R_{ij},
$$
where $r=\int_MRdV/\int_MdV$ is the average scalar curvature. The
factor $r$ serves to normalize the Ricci flow so that the volume
is constant. To see this we observe that $dV=\sqrt{\det
g_{ij}}\;dx$ and then
$$
\frac{\partial}{\partial t}\log\sqrt{\det g_{ij}}
=\frac{1}{2}g^{ij}\frac{\partial}{\partial t}g_{ij}=r-R,
$$
$$
\frac{d}{dt}\int_MdV=\int_M(r-R)dV=0.
$$
As observed by Hamilton \cite{Ha82}, the Ricci flow and the
normalized Ricci flow differ only by a change of scale in space
and a change of parametrization in time. Indeed, we first assume
that $g_{ij}(t)$ evolves by the (unnormalized) Ricci flow and
choose the normalization factor $\psi=\psi(t)$ so that
$\tilde{g}_{ij}=\psi g_{ij}$, and $\int_Md\tilde{\mu}=1$. Next we
choose a new time scale $\tilde{t}=\int\psi(t)dt$. Then for the
normalized metric $\tilde{g}_{ij}$ we have
$$
\tilde{R}_{ij}=R_{ij}, \tilde{R}=\frac{1}{\psi}R,
\tilde{r}=\frac{1}{\psi}r.
$$
Because $\int_Md\tilde{V}=1$, we see that
$\int_MdV=\psi^{-\frac{n}{2}}$. Then
\begin{align*}
\frac{d}{dt}\log\psi&  =\left(-\frac{2}{n}\right)\frac{d}{dt}\log\int_MdV\\
&  = \left(-\frac{2}{n}\right)\frac{\int_M\frac{\partial}{\partial
t}\sqrt{\det g_{ij}}\;dx}{\int_MdV}\\
&  = \frac{2}{n}r,
\end{align*}
since $\frac{\partial}{\partial t}g_{ij}=-2R_{ij}$ for the Ricci
flow. Hence it follows that
\begin{align*}
\frac{\partial}{\partial \tilde{t}}\tilde{g}_{ij}& =
\frac{\partial}{\partial t}g_{ij}+\left(\frac{d}{dt}\log\psi\right)g_{ij}\\
&  = \frac{2}{n}\tilde{r}\tilde{g}_{ij}-2\tilde{R}_{ij}.
\end{align*}
Thus studying the behavior of the Ricci flow near the maximal time
is equivalent to studying the long-time behavior of the normalized
Ricci flow.

In this chapter we will discuss long-time behavior of the normalized
Ricci flow for the following special cases: (1) compact
two-manifolds (cf. Hamilton \cite{Ha88} and Chow \cite{Chow}); (2)
compact three-manifolds with nonnegative Ricci curvature (cf.
Hamilton \cite{Ha82}); (3) compact four-manifolds with nonnegative
curvature operator (cf. Hamilton \cite{Ha86}); and (4) compact
three-manifolds with uniformly bounded normalized curvature (cf.
Hamilton \cite{Ha99}).


\section{The Ricci Flow on Two-manifolds}

Let $M$ be a compact surface, we will discuss in this section the
evolution of a Riemannian metric $g_{ij}$ under the normalized
Ricci flow. Most of the presentation in this section follows
Hamilton \cite{Ha88}, as well as Chow \cite{Chow}. We also refer
the reader to chapter 5 of Chow-Knopf's book \cite{CK04} for an
excellent description of the subject.

On a surface, the Ricci curvature is given by
$$
R_{ij}=\frac{1}{2}Rg_{ij}
$$
so the normalized Ricci flow equation becomes \begin{equation}
\frac{\partial}{\partial t}g_{ij}=(r-R)g_{ij}. %\eqno (5.1.1)
\end{equation} Recall the Gauss-Bonnet formula says
$$
\int_MRdV=4\pi\chi(M),
$$
where $\chi(M)$ is the Euler characteristic number of $M$. Thus
the average scalar curvature $r=4\pi\chi(M)/\int_MdV$ is constant
in time.

To obtain the evolution equation of the normalized curvature, we
recall a simple principle in \cite {Ha82} for converting from the
unnormalized to the normalized evolution equation on an
$n$-dimensional manifold. Let $P$ and $Q$ be two expressions
formed from the metric and curvature tensors, and let $\tilde{P}$
and $\tilde{Q}$ be the corresponding expressions for the
normalized Ricci flow. Since they differ by dilations, they differ
by a power of the normalized factor $\psi=\psi(t)$. We say $P$ has
\textbf{degree}\index{degree} $k$ if $\tilde{P}=\psi^kP$. Thus
$g_{ij}$ has degree 1, $R_{ij}$ has degree 0, $R$ has degree $-1$.

%\vskip 0.2cm \noindent {\bf Lemma 5.1.1} \;\emph{

%CZ-SOURCE-LINE-9936
\begin{lemma} [normalized evolution of homogeneous quantities]
  \label{lem:cz-normalized-evolution-homogeneous-quantities}
  \dcref{Cao--Zhu Lemma 5.1.1}
  \group{Normalized Ricci Flow on Surfaces}
  \source{cao-zhu}

Suppose $P$ satisfies
$$
\frac{\partial P}{\partial t}=\Delta P+Q
$$
for the unnormalized Ricci flow, and $P$ has degree $k$. Then $Q$
has degree $k-1$, and for the normalized Ricci flow,
$$
\frac{\partial \tilde{P}}{\partial \tilde{t}}=\tilde{\Delta}
\tilde{P}+\tilde{Q}+\frac{2}{n}k\tilde{r}\tilde{P}.
$$
\end{lemma}

%\vskip 0.1cm \noindent{\bf Proof.}
\begin{proof}
We first see $Q$ has degree $k-1$ since $\partial
\tilde{t}/\partial{t}=\psi$ and $\Delta=\psi\tilde{\Delta}$. Then
$$
\psi\frac{\partial}{\partial \tilde{t}}(\psi^{-k}\tilde{P})
=\psi\tilde{\Delta}(\psi^{-k}\tilde{P})+\psi^{-k+1}\tilde{Q}
$$
which implies
\begin{align*}
\frac{\partial \tilde{P}}{\partial \tilde{t}}&  =\tilde{\Delta}
\tilde{P}+\tilde{Q}+\frac{k}{\psi}\frac{\partial \psi}{\partial
\tilde{t}}\tilde{P}\\
&  = \tilde{\Delta}
\tilde{P}+\tilde{Q}+\frac{2}{n}k\tilde{r}\tilde{P}
\end{align*}
since $\frac{\partial}{\partial\tilde{t}}\log\psi
=(\frac{\partial}{\partial t}\log\psi)\psi^{-1}
=\frac{2}{n}\tilde{r}.$
%$$\eqno \#$$ \vskip 0.3cm
\end{proof}

We now come back to the normalized Ricci flow (5.1.1) on a compact
surface. By applying the above lemma to the evolution equation of
unnormalized scalar curvature, we have \begin{equation}
\frac{\partial R}{\partial t}=\Delta R+R^2-rR %\eqno (5.1.2)
\end{equation} for the normalized scalar curvature $R$. As a direct
consequence, by using the maximum principle, both nonnegative
scalar curvature and nonpositive scalar curvature are preserved
for the normalized Ricci flow on surfaces.

Let us introduce a potential function $\varphi$ as in the
K\"ahler-Ricci flow (see for example \cite{Cao85}). Since $R-r$
has mean value zero on a compact surface, there exists a unique
function $\varphi$, with mean value zero, such that \begin{equation}
\Delta \varphi=R-r.     %\eqno(5.1.3)
\end{equation} Differentiating (5.1.3) in time, we have
\begin{align*}
\frac{\partial }{\partial t}R
&  =\frac{\partial }{\partial t}(\Delta\varphi)\\
&  =(R-r)\Delta\varphi+g^{ij}\frac{\partial}{\partial
t}\left(\frac{\partial^2\varphi}{\partial x^i\partial x^j}
-\Gamma^k_{ij}\frac{\partial \varphi}{\partial x^k}\right)\\[4mm]
&  =(R-r)\Delta \varphi +\Delta\left(\frac{\partial \varphi}{\partial
t}\right).
\end{align*}
Combining with the equation (5.1.2), we get
$$
\Delta\left(\frac{\partial \varphi}{\partial t}\right)
=\Delta(\Delta\varphi)+r\Delta\varphi
$$
which implies that \begin{equation} \frac{\partial \varphi}{\partial
t}=\Delta\varphi+r\varphi-b(t)
%\eqno(5.1.4)
\end{equation} for some function $b(t)$ of time only. Since $\int_M\varphi
dV=0$ for all $t$, we have
\begin{align*}
0=\frac{d}{dt}\int_M\varphi d\mu&  =\int_M(\Delta
\varphi+r\varphi-b(t))d\mu+\int_M\varphi(r-R)d\mu \\
&  = -b(t)\int_Md\mu+\int_M|\nabla \varphi|^2d\mu.
\end{align*}
Thus the function $b(t)$ is given by
$$
b(t)=\frac{\int_M|\nabla\varphi|^2d\mu}{\int_Md\mu}.
$$

Define a function $h$ by
$$
h=\Delta\varphi+|\nabla\varphi|^2=(R-r)+|\nabla\varphi|^2,
$$
and set
$$
M_{ij}=\nabla_i\nabla_{j}\; \varphi-\frac{1}{2}\Delta\;\varphi
g_{ij}
$$
to be the traceless part of $\nabla_i\nabla_j\; \varphi$.

%\vskip 0.2cm \noindent {\bf Lemma 5.1.2} \;\emph{
%CZ-SOURCE-LINE-10027
\begin{lemma} [evolution of the curvature-potential energy density]
  \label{lem:cz-curvature-potential-density-evolution}
  \dcref{Cao--Zhu Lemma 5.1.2}
  \group{Normalized Ricci Flow on Surfaces}
  \source{cao-zhu}

The function $h$ satisfies the evolution equation \begin{equation}
\frac{\partial h}{\partial t}=\Delta h-2|M_{ij}|^2+rh.
%\eqno(5.1.5)
\end{equation}
\end{lemma}

%\vskip 0.1cm \noindent{\bf Proof.} \;
\begin{proof}
Under the normalized Ricci flow,
\begin{align*}
\frac{\partial}{\partial t}|\nabla \varphi|^2& =
\left(\frac{\partial}{\partial t}g^{ij}\right)
\nabla_i\varphi\nabla_j\varphi+2g^{ij}\left(\frac{\partial}{\partial
t}\nabla_i\varphi\right)(\nabla_j\varphi)\\
&  =(R-r)|\nabla\varphi|^2+2g^{ij}\nabla_i
(\Delta\varphi+r\varphi-b(t))\nabla_j\varphi\\
&  =(R+r)|\nabla\varphi|^2+2g^{ij}
(\Delta\nabla_i\varphi-R_{ik}\nabla_k\varphi)\nabla_j\varphi\\
&  =(R+r)|\nabla\varphi|^2+\Delta|\nabla\varphi|^2
-2|\nabla^2\varphi|^2-2g^{ij}R_{ik}\nabla_k\varphi\nabla_j\varphi\\
&
=\Delta|\nabla\varphi|^2-2|\nabla^2\varphi|^2+r|\nabla\varphi|^2,
\end{align*}
where $R_{ik}=\frac{1}{2}Rg_{ik}$ on a surface.

On the other hand we may rewrite the evolution equation (5.1.2) as
$$
\frac{\partial}{\partial t}(R-r)=\Delta(R-r)+(\Delta
\varphi)^2+r(R-r).
$$
Then the combination of above two equations yields
\begin{align*}
\frac{\partial}{\partial t}h&  =\Delta
h-2(|\nabla^2\varphi|^2-\frac{1}{2}(\Delta\varphi)^2)+rh\\
&  =\Delta h-2|M_{ij}|^2+rh
\end{align*}
as desired.
%$$\eqno \#$$ \vskip 0.3cm
\end{proof}

As a direct consequence of the evolution equation (5.1.5) and the
maximum principle, we have \begin{equation}
R\leq C_1e^{rt}+r %\eqno(5.1.6)
\end{equation} for some positive constant $C_1$ depending only on the initial
metric.

On the other hand, it follows from (5.1.2) that
$R_{\min}(t)=\min_{x\in M}R(x,t)$ satisfies
$$
\frac{d}{dt}R_{\min}\geq R_{\min}(R_{\min}-r)\geq 0
$$
whenever $R_{\min}\leq 0$. This says that \begin{equation}
R_{\min}(t)\geq-C_2, \quad\text{ for all }\; t>0 %\eqno (5.1.7)
\end{equation} for some positive constant $C_2$ depending only on the initial
metric.

Thus the combination of (5.1.6) and (5.1.7) implies the following
long time existence result.

%\vskip 0.2cm \noindent {\bf Proposition 5.1.3} \;\emph{
%CZ-SOURCE-LINE-10088
\begin{proposition} [immortality of normalized surface Ricci flow]
  \label{prop:cz-normalized-surface-flow-immortal}
  \dcref{Cao--Zhu Proposition 5.1.3}
  \group{Normalized Ricci Flow on Surfaces}
  \source{cao-zhu}

For any initial metric on a compact surface, the normalized Ricci
flow $(5.1.1)$ has a solution for all time.
%$$\eqno \#$$
\end{proposition}

To investigate the long-time behavior of the solution, let us now
divide the discussion into three cases: $\chi(M)<0$; $\chi(M)=0$;
and $\chi(M)>0$.

\medskip
{\it Case} (1): $\chi(M)<0$ (i.e., $r<0$).

\smallskip
{}From the evolution equation (5.1.2), we have
\begin{align*}
\frac{d}{dt}R_{\min}&  \geq R_{\min}(R_{\min}-r) \\
&  \geq r(R_{\min}-r), \; \text{ on } \; M\times [0,+\infty)
\end{align*}
which implies that
$$
R-r\geq-\tilde{C_1}e^{rt}, \; \text{ on } \; M\times[0,+\infty)
$$
for some positive constant $\tilde{C_1}$ depending only on the
initial metric. Thus by combining with (5.1.6) we have \begin{equation}
-\tilde{C_1}e^{rt}\leq R-r \leq C_1e^{rt}, \; \text{ on } \;
M\times
[0,+\infty). %\eqno (5.1.8)
\end{equation}

%\vskip 0.2cm \noindent {\bf Theorem 5.1.4} (
%CZ-SOURCE-LINE-10119
\begin{theorem}[convergence on surfaces of negative Euler characteristic]
  \label{thm:cz-negative-euler-surface-convergence}
  \dcref{Cao--Zhu Theorem 5.1.4}
  \group{Normalized Ricci Flow on Surfaces}
  \source{cao-zhu}

On a compact surface with $\chi(M)<0$, for any initial metric the
solution of the normalized Ricci flow $(5.1.1)$ exists for all
time and converges in the $C^{\infty}$ topology to a metric with
negative constant curvature.
\end{theorem}

%\vskip 0.1cm \noindent {\bf Proof.}
\begin{proof}
The estimate (5.1.8) shows that the scalar curvature $R(x,t)$
converges exponentially to the negative constant $r$ as
$t\rightarrow+\infty$.

Fix a tangent vector $v\in T_xM$ at a point $x\in M$ and let
$|v|^2_t=g_{ij}(x,t)v^iv^j$. Then we have
\begin{align*}
\frac{d}{dt}|v|_t^2
&  =\left(\frac{\partial}{\partial t}g_{ij}(x,t)\right)v^iv^j\\
&  = (r-R)|v|_t^2
\end{align*}
which implies
$$
\left|\frac{d}{dt}\log|v|_t^2\right| \leq Ce^{rt}, \; \text{ for
all } \; t>0
$$
for some positive constant $C$ depending only on the initial
metric (by using (5.1.8)). Thus $|v|_t^2$ converges uniformly to a
continuous function $|v|_{\infty}^2$ as $t\rightarrow +\infty$ and
$|v|_{\infty}^2\neq 0$ if $v\neq 0$. Since the parallelogram law
continues to hold to the limit, the limiting norm $|v|_{\infty}^2$
comes from an inner product $g_{ij}(\infty)$. This says, the
metrics $g_{ij}(t)$ are all equivalent and as $t\rightarrow
+\infty$, the metric $g_{ij}(t)$ converges uniformly to a
positive-definite metric tensor $g_{ij}(\infty)$ which is
continuous and equivalent to the initial metric.

By the virtue of Shi's derivative estimates of the unnormalized
Ricci flow in Section 1.4, we see that all derivatives and higher
order derivatives of the curvature of the solution $g_{ij}$ of the
normalized flow are uniformly bounded on $M\times [0,+\infty)$.
This shows that the limiting metric $g_{ij}(\infty)$ is a smooth
metric with negative constant curvature and the solution
$g_{ij}(t)$ converges to the limiting metric $g_{ij}(\infty)$ in
the $C^{\infty}$ topology as $t\rightarrow +\infty$.
%$$\eqno\#$$
\end{proof}

\smallskip
{\it Case} (2): $\chi(M)=0$, i.e., $r=0$.

\smallskip
{}The following argument of dealing with the case $\chi(M)=0$ is
adapted from Chow-Knopf's book (cf. section 5.6 of \cite{CK04}).
From (5.1.6) and (5.1.7) we know that the curvature remains
bounded above and below. To get the convergence, we consider the
potential function $\varphi$ of (5.1.3) again. The evolution of
$\varphi$ is given by (5.1.4). We renormalize the function
$\varphi$ by
$$
\tilde{\varphi}(x,t) =\varphi(x,t)+\int b(t)dt, \qquad \text{on
}\;M\times[0,+\infty).
$$
Then, since $r=0$, $\tilde{\varphi}$ evolves by \begin{equation} \frac{\partial
\tilde{\varphi}}{\partial t} =\Delta\tilde{\varphi},\qquad
\text{on}
\;M\times[0,+\infty).  %\eqno(5.1.9)
\end{equation} {}From the proof of Lemma 5.1.2, we get \begin{equation} \frac{\partial
}{\partial t}|\nabla \tilde{\varphi}|^2= \Delta|\nabla
\tilde{\varphi}|^2-2|\nabla^2\tilde{\varphi}|^2. %\eqno(5.1.10)
\end{equation} Clearly, we have \begin{equation} \frac{\partial }{\partial
t}\tilde{\varphi}^2
=\Delta\tilde{\varphi}^2-2|\nabla\tilde{\varphi}|^2. %\eqno(5.1.11)
\end{equation} Thus it follows that
$$
\frac{\partial }{\partial t}(t|\nabla\tilde{\varphi}|^2
+\tilde{\varphi}^2)\le\Delta(t|\nabla \tilde{\varphi}|^2+
\tilde{\varphi}^2).
$$
Hence by applying the maximum principle, there exists a positive
constant $C_3$ depending only on the initial metric such that \begin{equation}
|\nabla\tilde{\varphi}|^2(x,t)\le \frac{C_3}{1+t},
\qquad \text{on}\;M\times[0,+\infty). %\eqno(5.1.12)
\end{equation} In the following we will use this decay estimate to obtain a
decay estimate for the scalar curvature.

By the evolution equations (5.1.2) and (5.1.10), we have
\begin{align*}
\frac{\partial }{\partial t}(R+2|\nabla\tilde{\varphi}|^2) &
=\Delta(R+2|\nabla\tilde{\varphi}|^2)
+R^2-4|\nabla^2\tilde{\varphi}|^2\\
&  \le \Delta(R+2|\nabla\tilde{\varphi}|^2)-R^2
\end{align*}
since $R^2=(\Delta\tilde{\varphi})^2\le
2|\nabla^2\tilde{\varphi}|^2$. Thus by using (5.1.12), we have
\begin{align*}
& \frac{\partial }{\partial
t}[t(R+2|\nabla\tilde{\varphi}|^2)]\\
&  \le \Delta[t(R+2|\nabla\tilde{\varphi}|^2)]
-tR^2+R+2|\nabla\tilde{\varphi}|^2\\
& \leq\Delta[t(R+2|\nabla\tilde{\varphi}|^2)]
-t(R+2|\nabla\tilde{\varphi}|^2)^2+(1+4t|\nabla
\tilde{\varphi}|^2)(R+2|\nabla\tilde{\varphi}|^2)\\
&  \le \Delta[t(R+2|\nabla\tilde{\varphi}|^2)]
-[t(R+2|\nabla\tilde{\varphi}|^2)
-(1+4C_3)](R+2|\nabla\tilde{\varphi}|^2)\\
&  \le \Delta[t(R+2|\nabla\tilde{\varphi}|^2)]
\end{align*}
wherever $t(R+2|\nabla\tilde{\varphi}|^2)\ge (1+4C_3)$. Hence by
the maximum principle, there holds \begin{equation}
R+2|\nabla\tilde{\varphi}|^2\le\frac{C_4}{1+t},\qquad
\text{on}\;M\times[0,+\infty)  %\eqno(5.1.13)
\end{equation} for some positive constant $C_4$ depending only on the initial
metric.

On the other hand, the scalar curvature satisfies
$$
\frac{\partial R}{\partial t} =\Delta R+R^2,\qquad \text{on
}\;M\times[0,+\infty).
$$
It is not hard to see that \begin{equation}
R\ge\frac{R_{\min}(0)}{1-R_{\min}(0)t},\qquad
\text{on }\;M\times[0,+\infty),  %NEW, missing number
\end{equation} by using the maximum principle. So we obtain the decay
estimate for the scalar curvature \begin{equation}
|R(x,t)|\le\frac{C_5}{1+t},\qquad \text{on }\;M\times[0,+\infty),
%\eqno(5.1.15)
\end{equation} for some positive constant $C_5$ depending only on the initial
metric.

%\vskip 0.2cm \noindent {\bf Theorem 5.1.5} (
%CZ-SOURCE-LINE-10249
\begin{theorem}[convergence on surfaces of zero Euler characteristic]
  \label{thm:cz-zero-euler-surface-convergence}
  \dcref{Cao--Zhu Theorem 5.1.5}
  \group{Normalized Ricci Flow on Surfaces}
  \source{cao-zhu}

On a compact surface with $\chi(M)=0$, for any initial metric the
solution of the normalized Ricci flow $(5.1.1)$ exists for all
time and converges in $C^\infty$ topology to a flat metric.
\end{theorem}

%\vskip 0.1cm \noindent{\bf Proof.}
\begin{proof}
Since $\frac{\partial \tilde{\varphi}}{\partial
t}=\Delta\tilde{\varphi}$, it follows from the maximum principle
that
$$
|\tilde\varphi(x,t)|\le C_6,\qquad \text{on }\; M\times
[0,+\infty)
$$
for some positive constant $C_6$ depending only on the initial
metric. Recall that $\Delta\tilde{\varphi}=R$. We thus obtain for
any tangent vector $v\in T_xM$ at a point $x\in M$,
\begin{align*}
\frac{d}{dt}|v|^2_t&
= \left(\frac{\partial }{\partial t}g_{ij}(x,t)\right)v^iv^j\\
&  = -R(x,t)|v|^2_t
\end{align*}
and then
\begin{align*}
\left|\log\frac{|v|^2_t}{|v|^2_0}\right|
&  = \left|\int_0^t\frac{d}{dt}\log\right|v_t|^2dt|\\
&  = \left|\int_0^tR(x,t)dt\right|\\
&  = \left|\tilde{\varphi}(x,t)-\tilde{\varphi}(x,0)\right|\\
&  \le 2C_6,
\end{align*}
for all $x\in M$ and $t\in[0,+\infty).$ This shows that the
solution $g_{ij}(t)$ of the normalized Ricci flow are all
equivalent. This gives us control of the diameter and injectivity
radius.

As before, by Shi's derivative estimates of the unnormalized Ricci
flow, all derivatives and higher order derivatives of the
curvature of the solution $g_{ij}$ of the normalized Ricci flow
(5.1.1) are uniformly bounded on $M\times[0,+\infty)$. By the
virtue of Hamilton's compactness theorem (Theorem 4.1.5) we see
that the solution $g_{ij}(t)$ subsequentially converges in
$C^\infty$ topology. The decay estimate (5.1.15) implies that each
limit must be a flat metric on $M$. Clearly, we will finish the
proof if we can show that limit is unique.

Note that the solution $g_{ij}(t)$ is changing conformally under
the Ricci flow (5.1.1) on surfaces. Thus each limit must be
conformal to the initial metric, denoted by $\bar{g}_{ij}$. Let us
denote $g_{ij}(\infty)=e^u\bar{g}_{ij}$ to be a limiting metric.
Since $g_{ij}(\infty)$ is flat, it is easy to compute
$$
0=e^{-u}(\bar{R}-\bar{\Delta}u),\qquad \mbox{on }\;M,
$$
where $\bar{R}$ is the curvature of $\bar{g}_{ij}$ and
$\bar{\Delta}$ is the Laplacian in the metric $\bar{g}_{ij}$. The
solution of Poission equation
$$\bar{\Delta}u=\bar{R},\qquad \mbox{on }\;M$$
is unique up to constant. Moreover the constant must be also
uniquely determined since the area of the solution of the
normalized Ricci flow (5.1.1) is constant in time. So the limit is
unique and we complete the proof of Theorem 5.1.5.
%$$\eqno\#$$\vskip 0.3cm
\end{proof}

\medskip
{\it Case} (3): $\chi(M)>0$, i.e., $r>0$.

\smallskip
This is the most difficult case. The first proof is due to Ben
Chow\cite{Chow}, based on the important work of Hamilton
\cite{Ha88}.  By now there exist several proofs (cf.
Bartz-Struwe-Ye \cite{BSY} and Struwe \cite{St}, etc). But, in
contrast to the previous two cases, none of the proofs depends
only on maximum principle type argument. In fact, all the proofs
rely on some kind of combination of the maximum principle argument
and certain integral estimate of the curvature.

In the pioneer work \cite{Ha88}, Hamilton considered the important
special case when the initial metric is of positive scalar
curvature. He introduced an integral quantity
$$
E=\int_M R \log R \;dV,
$$
which he called {\bf entropy}, for the (normalized) Ricci flow on
a surface $M$ with positive curvature, and showed that the entropy
is monotone decreasing under the flow. By combining this entropy
estimate with the Harnack inequality for the curvature (Corollary
2.5.3), Hamilton obtained the uniform bound on the curvature of
the normalized Ricci flow on $M$ with positive curvature.
Furthermore, he showed that the evolving metric converges to a
shrinking Ricci soliton on $M$ and that the shrinking Ricci
soliton must be a round metric on the $2$-sphere $\mathbb S^2$.
Subsequently, Chow \cite{Chow} extended Hamilton's work to the
general case when the curvature may change signs. More precisely,
he proved that given any initial metric on a compact surface $M$
with $\chi(M)>0$, the evolving metric under the (normalized) Ricci
flow will have positive curvature after a finite time. Hence, when
combined with Hamilton's result, B. Chow's result implies that the
solution under the normalized flow on $M$ converges to the round
metric on $\mathbb S^2$.

In the following we will basically follow the arguments of
Hamilton \cite{Ha88} and Chow \cite{Chow}, except when we prove
the uniform bound of the evolving scalar curvature we will present
a new argument using the Harnack inequality of Chow \cite{Chow}
and Perelman's no local collapsing theorem I$'$ (as was done in
the joint work of Bing-Long Chen and the authors \cite{CCZ} where
they considered the K\"ahler-Ricci flow of nonnegative holomorphic
bisectional curvature).

Given any initial metric on $M$ with $\chi(M)>0$, we consider the
solution $g_{ij}(t)$ of the normalized Ricci flow (5.1.1). Recall
that the (scalar) curvature $R$ satisfies the evolution equation
$$
\frac{\partial }{\partial t}R=\Delta R+R^2-rR.
$$
The corresponding ODE is \begin{equation}
\frac{ds}{dt}=s^2-rs.  %\eqno(5.1.16)
\end{equation}

Let us choose $c>1$ and close to 1 so that $r/(1-c)<\min_{x\in
M}R(x,0)$. It is clear that the function $s(t)=r/(1-ce^{rt})<0$ is
a solution of the ODE (5.1.16) with $s(0)<\min\limits_{x\in
M}R(x,0)$. Then the difference of $R$ and $s$ evolves by \begin{equation}
\frac{\partial }{\partial t}(R-s)=\Delta(R-s)+(R-r+s)(R-s).
%\eqno(5.1.17)
\end{equation} Since $\min_{x\in M}R(x,0)-s(0)>0$, the maximum principle
implies that $R-s>0$ for all times.

First, we need the Harnack inequality obtained by B. Chow
\cite{Chow}, which is an extension of Theorem 2.5.2, for the
normalized Ricci flow whose curvature may change signs.

Consider the quantity
$$
L=\log(R-s).
$$
It is easy to compute
$$
\frac{\partial L}{\partial t}=\Delta L+|\nabla L|^2+R-r+s.
$$
Set
$$
Q=\frac{\partial L}{\partial t}-|\nabla L|^2-s=\Delta L+R-r.
$$
By a direct computation and using the estimate (5.1.8), we have
\begin{align*}
\frac{\partial }{\partial t}Q&  = \Delta\left(\frac{\partial
L}{\partial t}\right)
+(R-r)\Delta L+\frac{\partial R}{\partial t}\\
&  =\Delta Q+2|\nabla^2L|^2+2\langle\nabla L,\nabla(\Delta
L)\rangle+R|\nabla L|^2\\
&\quad +(R-r)\Delta L+\Delta R+R(R-r)\\
&  =\Delta Q+2|\nabla^2L|^2+2\langle\nabla L,\nabla
Q\rangle+2(R-r)\Delta L+(R-r)^2\\
&\quad +(r-s)\Delta L+s|\nabla L|^2+r(R-r)\\
&  = \Delta Q+2\langle\nabla L,\nabla Q\rangle
+2|\nabla^2L|^2+2(R-r)\Delta L+(R-r)^2\\
&\quad +(r-s)Q+s|\nabla L|^2+s(R-r)\\
&  \ge \Delta Q+2\langle \nabla L,\nabla
Q\rangle+Q^2+(r-s)Q+s|\nabla L|^2-C.
\end{align*}
Here and below $C$ is denoted by various positive constants
depending only on the initial metric.

In order to control the bad term $s|\nabla L|^2$, we consider
\begin{align*}
\frac{\partial }{\partial t}(sL)
&  =\Delta(sL)+s|\nabla L|^2+s(R-r+s)+s(s-r)L\\
&  \ge \Delta(sL)+2\langle\nabla L,\nabla(sL)\rangle -s|\nabla
L|^2-C
\end{align*}
by using the estimate (5.1.8) again. Thus
\begin{align*}
\frac{\partial }{\partial t}(Q+sL) &
\ge\Delta(Q+sL)+2\langle\nabla
L,\nabla(Q+sL)\rangle+Q^2+(r-s)Q-C\\
&  \ge \Delta(Q+sL)+2\langle\nabla L,\nabla(Q+sL)\rangle
+\frac{1}{2}[(Q+sL)^2-C^2],
\end{align*}
since $sL$ is bounded by (5.1.8). This, by the maximum principle,
implies that
$$
Q\ge -C,\qquad \text{for\;all }\;t\in[0,+\infty).
$$
Then for any two points $x_1,x_2\in M$ and two times $t_2>t_1\ge
0$, and a path $\gamma:[t_1,t_2]\rightarrow M$ connecting $x_1$ to
$x_2$, we have
\begin{align*}
L(x_2,t_2)-L(x_1,t_1)
&  = \int_{t_1}^{t_2}\frac{d}{dt}L(\gamma(t),t)dt\\
&  = \int_{t_1}^{t_2}\left(\frac{\partial L}{\partial t}+\langle\nabla
L,\dot{\gamma}\rangle\right)dt\\
&  \ge -\frac{1}{4}\Delta-C(t_2-t_1)
\end{align*}
where
\begin{align*}
\Delta&  = \Delta(x_1,t_1;x_2,t_2)\\
&  =
\inf\left\{\int_{t_1}^{t_2}|\dot{\gamma}(t)|^2_{g_{ij}(t)}dt\, |\,
\gamma:[t_1,t_2]\rightarrow
M\;\text{with}\;\gamma(t_1)=x_1,\gamma(t_2)=x_2\right\}.
\end{align*}
Thus we have proved the following Harnack inequality of B. Chow.

%\vskip 0.2cm \noindent{\bf Lemma 5.1.6} (
%CZ-SOURCE-LINE-10456
\begin{lemma}[space-time Harnack estimate on positive-Euler surfaces]
  \label{lem:cz-positive-euler-surface-harnack}
  \dcref{Cao--Zhu Lemma 5.1.6}
  \group{Normalized Ricci Flow on Surfaces}
  \source{cao-zhu}

There exists a positive constant C depending only on the initial
metric such that for any $x_1$, $x_2\in M$ and $t_2>t_1\ge 0$,
$$
R(x_1,t_1)-s(t_1) \le
e^{\frac{\Delta}{4}+C(t_2-t_1)}(R(x_2,t_2)-s(t_2))
$$
where
$$
\Delta=\inf\left\{\int_{t_1}^{t_2}|\dot{\gamma}(t)|^2_tdt\ |\
\gamma:[t_1,t_2]\rightarrow M\;with
\;\gamma(t_1)=x_1,\gamma(t_2)=x_2\right\}.
$$
%$$\eqno \#$$ \vskip 0.3cm
\end{lemma}

We now state and prove the following uniform bound estimate for
the curvature, a consequence of the results of Hamilton
\cite{Ha88} and Chow \cite{Chow}. We remark the special case when
the scalar curvature $R>0$ is first proved by Hamilton
\cite{Ha88}. As we mentioned before, the proof here is adapted
from \cite{CCZ}.

%\vskip 0.2cm \noindent {\bf Proposition 5.1.7} \;\emph{

%CZ-SOURCE-LINE-10481
\begin{proposition} [uniform scalar-curvature bounds on positive-Euler surfaces]
  \label{prop:cz-positive-euler-surface-curvature-bounds}
  \dcref{Cao--Zhu Proposition 5.1.7}
  \group{Normalized Ricci Flow on Surfaces}
  \source{cao-zhu}

Let $(M,g_{ij}(t))$ be a solution of the normalized Ricci flow on
a compact surface with $\chi(M)>0$. Then there exist a time
$t_0>0$ and a positive constant $C$ such that the estimate
$$
C^{-1}\le R(x,t)\le C
$$
holds for all $x\in M$ and $t\in [t_0,+\infty).$
\end{proposition}

%\vskip 0.1cm\noindent{\bf Proof.} \;
\begin{proof}
  \uses{lem:cz-positive-euler-surface-harnack}
Recall that
$$
R(x,t)\ge s(t) =\frac{r}{1-ce^{rt}},\qquad \text{on}\; M\times
[0,+\infty).
$$
For any $\varepsilon\in(0,r)$, there exists a large enough $t_0>0$
such that \begin{equation} R(x,t)\ge-\varepsilon^2,\qquad \text{on}\;
M\times[t_0,+\infty).
%\eqno(5.1.18)
\end{equation} Let $t$ be any fixed time with $t\ge t_0+1$. Obviously there
is some point $x_0\in M$ such that $R(x_0,t+1)=r$.

Consider the geodesic ball $B_t(x_0,1)$, centered at $x_0$ and
radius 1 with respect to the metric at the fixed time $t$. For any
point $x\in B_t(x_0,1)$, we choose a geodesic $\gamma$:
$[t,t+1]\rightarrow M$ connecting $x$ and $x_0$ with respect to
the metric at the fixed time $t$. Since
$$
\frac{\partial }{\partial t}g_{ij}=(r-R)g_{ij} \le 2rg_{ij}\qquad
\text{on }\;M\times[t_0,+\infty),
$$
we have
$$
\int_t^{t+1}|\dot{\gamma}(\tau)|^2_\tau d\tau\le
e^{2r}\int_t^{t+1} |\dot{\gamma}(\tau)|^2_t d\tau\leq e^{2r}.
$$
Then by Lemma 5.1.6, we have
\begin{align}
R(x,t)& \leq s(t)+\exp\left\{\frac{1}{4}e^{2r}+C\right\}
\cdot(R(x_0,t+1)-s(t+1))\\
& \leq C_1,\qquad \text{as }\; x\in B_t(x_0,1),\nonumber
\end{align} %\eqno(5.1.19)
for some positive constant $C_1$ depending only on the initial
metric. Note the the corresponding unnormalized Ricci flow in this
case has finite maximal time since its volume decreases at a fixed
rate $-4\pi \chi (M) < 0$. Hence the no local collapsing theorem
I$'$ (Theorem 3.3.3) implies that the volume of $B_t(x_0,1)$ with
respect to the metric at the fixed time $t$ is bounded from below
by \begin{equation}
\Vol_t(B_t(x_0,1))\geq C_2 %\eqno(5.1.20)
\end{equation} for some positive constant $C_2$ depending only on the initial
metric.

We now want to bound the diameter of $(M,g_{ij}(t))$ from above.
The following argument is analogous to Yau in \cite{Y76} where he
got a lower bound for the volume of geodesic balls of a complete
Riemannian manifold with nonnegative Ricci curvature. Without loss
of generality, we may assume that the diameter of $(M,g_{ij}(t))$
is at least 3. Choose a point $x_1\in M$ such that the distance
$d_t(x_0,x_1)$ between $x_1$ and $x_0$ with respect to the metric
at the fixed time $t$ is at least a half of the diameter of
$(M,g_{ij}(t))$. By (5.1.18), the standard Laplacian comparison
theorem (c.f. \cite{ScY}) implies
$$
\Delta \rho^2=2\rho\Delta\rho+2\leq2(1+\varepsilon\rho)+2
$$
in the sense of distribution, where $\rho$ is the distance
function from $x_1$ (with respect to the metric $g_{ij}(t))$. That
is, for any $\varphi\in C_0^\infty(M)$, $\varphi\geq0$, we have
\begin{equation} -\int_M\nabla\rho^2\cdot\nabla\varphi
\leq\int_M[2(1+\varepsilon\rho)+2]\varphi. %\eqno(5.1.21)
\end{equation} Since $C_0^\infty(M)$ functions can be approximated by
Lipschitz functions in the above inequality, we can set
$\varphi(x)=\psi(\rho(x))$, $x\in M$, where $\psi(s)$ is given by
$$
\psi(s)=\begin{cases}
1,& 0\leq s\leq d_t(x_0,x_1)-1,\\
\psi'(s)=-\frac{1}{2},& d_t(x_0,x_1)-1\leq s\leq
d_t(x_0,x_1)+1,\\
0,& s\geq d_t(x_0,x_1)+1.
\end{cases}
$$
Thus, by using (5.1.20), the left hand side of (5.1.21) is
\begin{displaymath}
\begin{split}
&-\int_M\nabla\rho^2\cdot\nabla\varphi \\
&  =\int_{B_t(x_1,d_t(x_0,x_1)+1)
\setminus B_t(x_1,d_t(x_0,x_1)-1)}\rho\\
&  \geq(d_t(x_0,x_1)-1)\Vol_t(B_t(x_1,d_t(x_0,x_1)+1)
\setminus B_t(x_1,d_t(x_0,x_1)-1))\\
&  \geq(d_t(x_0,x_1)-1)\Vol_t(B_t(x_0,1))\\
&  \geq(d_t(x_0,x_1)-1)C_2,
\end{split}
\end{displaymath}
 and the right hand side of (5.1.21) is
\begin{displaymath}
\begin{split}
\int_M[2(1+\varepsilon\rho)+2]\varphi
&  \leq\int_{B_t(x_1,d_t(x_0,x_1)+1)}[2(1+\varepsilon\rho)+2]\\
&  \leq [2(1+\varepsilon d_t(x_0,x_1))+4]
\Vol_t(B_t(x_1,d_t(x_0,x_1)+1))\\
&  \leq [2(1+\varepsilon d_t(x_0,x_1))+4]A
\end{split}
\end{displaymath}
where $A$ is the area of $M$ with respect to the initial metric.
Here we have used the fact that the area of solution of the
normalized Ricci flow is constant in time. Hence
$$
C_2(d_t(x_0,x_1)-1)\leq[2(1+\varepsilon d_t(x_0,x_1))+4] A,
$$
which implies, by choosing $\varepsilon>0$ small enough,
$$
d_t(x_0,x_1)\leq C_3
$$
for some positive constant $C_3$ depending only on the initial
metric. Therefore, the diameter of $(M,g_{ij}(t))$ is uniformly
bounded above by \begin{equation}
{\rm diam}\,(M,g_{ij}(t))\leq 2C_3  %\eqno(5.1.22)
\end{equation} for all $t\in[t_0,+\infty)$.

We then argue, as in deriving (5.1.19), by applying Lemma 5.1.6
again to obtain
$$
R(x,t)\leq C_4,\ \ \ \text{on }\; M\times[t_0,+\infty)
$$
for some positive constant $C_4$ depending only on the initial
metric.

It remains to prove a positive lower bound estimate of the
curvature. First, we note that the function $s(t)\to 0$ as
$t\rightarrow+\infty$, and the average scalar curvature of the
solution equals to $r$, a positive constant. Thus the Harnack
inequality in Lemma 5.1.6 and the diameter estimate (5.1.22) imply
a positive lower bound for the curvature. Therefore we have
completed the proof of Proposition 5.1.7.
\end{proof}
%$$\eqno \#$$ \vskip 0.3cm

Next we consider long-time convergence of the normalized flow.

Recall that the trace-free part of the Hessian of the potential
$\varphi$ of the curvature is the tensor $M_{ij}$ defined by
$$
M_{ij}=\nabla_i\nabla_j\varphi-\frac{1}{2}\Delta\varphi\cdot
g_{ij},
$$
where by (5.1.3),
$$
\Delta\varphi=R-r.
$$

%\vskip 0.2cm\noindent{\bf Lemma 5.1.8} \ \emph{
%CZ-SOURCE-LINE-10635
\begin{lemma} [decay evolution for the trace-free Hessian]
  \label{lem:cz-tracefree-hessian-evolution}
  \dcref{Cao--Zhu Lemma 5.1.8}
  \group{Normalized Ricci Flow on Surfaces}
  \source{cao-zhu}

We have \begin{equation} \frac{\partial}{\partial t}|M_{ij}|^2
=\Delta|M_{ij}|^2-2|\nabla_k M_{ij}|^2-2R|M_{ij}|^2,
\; \text{ on }\; M\times[0,+\infty). %\eqno(5.1.23)
\end{equation}
\end{lemma}

%\vskip 0.1cm\noindent{\bf Proof.}\
\begin{proof} This follows from a standard computation
(e.g., cf. Editors' note on p. 217 of \cite{CCCY}).

First we note the time-derivative of the Levi-Civita connection is
\begin{displaymath}
\begin{split}
\frac{\partial}{\partial t}\Gamma^k_{ij} &
=\frac{1}{2}g^{kl}\left(\nabla_i\frac{\partial}{\partial t}g_{jl}
+\nabla_j\frac{\partial}{\partial t}g_{il}
-\nabla_l\frac{\partial}{\partial t}g_{ij}\right)\\
&  =\frac{1}{2}\left(-\nabla_iR\cdot\delta_j^k
-\nabla_jR\cdot\delta_i^k+\nabla^kR\cdot g_{ij}\right).
\end{split}
\end{displaymath}
By using this and (5.1.4), we have
\begin{displaymath}
\begin{split}
  \frac{\partial}{\partial t}M_{ij}
& = \nabla_i\nabla_j\left(\frac{\partial \varphi}{\partial t}\right)
    -\left(\frac{\partial}{\partial t}\Gamma^k_{ij}\right)\nabla_k\varphi
    -\frac{1}{2}\frac{\partial}{\partial t}[(R-r)g_{ij}]\\
& = \nabla_i\nabla_j\Delta\varphi
+\frac{1}{2}(\nabla_iR\cdot\nabla_j\varphi
    +\nabla_jR\cdot\nabla_i\varphi
-\langle\nabla R,\nabla\varphi\rangle g_{ij})\\
& \quad -\frac{1}{2}\Delta R\cdot g_{ij}+rM_{ij}.
\end{split}
\end{displaymath}
Since on a surface,
$$
R_{ijkl}=\frac{1}{2}R(g_{il}g_{jk}-g_{ik}g_{jl}),
$$
we have
\begin{displaymath}
\begin{split}
& \nabla_i\nabla_j\Delta\varphi \\
& = \nabla_i\nabla_k\nabla_j\nabla^k\varphi
-\nabla_i(R_{jl}\nabla^l\varphi)\\
& = \nabla_k\nabla_i\nabla_j\nabla^k\varphi
-R^l_{ikj}\nabla_l\nabla^k\varphi-R_{il}\nabla_j\nabla^l\varphi
    -R_{jl}\nabla_i\nabla^l\varphi
-\nabla_iR_{jl}\nabla^l\varphi\\
& = \Delta\nabla_i\nabla_j\varphi
-\nabla^k(R^l_{ikj}\nabla_l\varphi)-R^l_{ikj}\nabla_l\nabla^k\varphi\\
&\quad  -R_{il}\nabla_j\nabla^l\varphi
-R_{jl}\nabla_i\nabla^l\varphi-\nabla_iR_{jl}\nabla^l\varphi\\
& = \Delta\nabla_i\nabla_j\varphi
-\frac{1}{2}(\nabla_iR\cdot\nabla_j\varphi+\nabla_jR\cdot\nabla_i\varphi
    -\langle\nabla R,\nabla \varphi\rangle g_{ij})\\
&\quad  -2R\left(\nabla_i\nabla_j\varphi-\frac{1}{2}\Delta\varphi\cdot
    g_{ij}\right).
\end{split}
\end{displaymath}
Combining these identities, we get
\begin{displaymath}
\begin{split}
\frac{\partial}{\partial t}M_{ij} & =
\Delta\nabla_i\nabla_j\varphi-\frac{1}{2}\Delta R\cdot g_{ij}
+(r-2R)M_{ij}\\
& = \Delta\left(\nabla_i\nabla_j\varphi-\frac{1}{2}(R-r)g_{ij}\right)
+(r-2R)M_{ij}.
\end{split}
\end{displaymath}
Thus the evolution $M_{ij}$ is given by \begin{equation} \frac{\partial
M_{ij}}{\partial t}=\Delta M_{ij}+(r-2R)M_{ij}.
%\eqno(5.1.24)
\end{equation} Now the lemma follows from (5.1.24) and a straightforward
computation.
%$$\eqno\#$$
\end{proof}

Proposition 5.1.7 tells us that the curvature $R$ of the solution
to the normalized Ricci flow is uniformly bounded from below by a
positive constant for $t$ large. Thus we can apply the maximum
principle to the equation (5.1.23) in Lemma 5.1.8 to obtain the
following estimate.

%\vskip 0.2cm \noindent {\bf Proposition 5.1.9} \emph{\
%CZ-SOURCE-LINE-10721
\begin{proposition}  [exponential decay of the trace-free Hessian]
  \label{prop:cz-tracefree-hessian-exponential-decay}
  \dcref{Cao--Zhu Proposition 5.1.9}
  \group{Normalized Ricci Flow on Surfaces}
  \source{cao-zhu}

Let $(M,g_{ij}(t))$ be a solution of the normalized Ricci flow on
a compact surface with $\chi(M)>0$. Then there exist positive
constants $c$ and $C$ depending only on the initial metric such
that
$$
|M_{ij}|^2\leq Ce^{-ct},\quad \text{ on }\; M\times[0,+\infty).
$$
%$$\eqno \#$$
\end{proposition}

Now we consider a modification of the normalized Ricci flow.
Consider the equation \begin{equation} \frac{\partial}{\partial
t}g_{ij}=2M_{ij}
=(r-R)g_{ij}+2\nabla_i\nabla_j\varphi. %\eqno(5.1.25)
\end{equation} As we saw in Section 1.3, the solution of this modified flow
differs from that of the normalized Ricci flow only by a one
parameter family of diffeomorphisms generated by the gradient
vector field of the potential function $\varphi$. Since the
quantity $|M_{ij}|^2$ is invariant under diffeomorphisms, the
estimate $|M_{ij}|^2\leq Ce^{-ct}$ also holds for the solution of
the modified flow (5.1.25). This exponential decay estimate then
implies the solution $g_{ij}(x,t)$ of the modified flow (5.1.25)
converges exponentially to a continuous metric $g_{ij}(\infty)$ as
$t\rightarrow+\infty$. Furthermore, by the virtue of Hamilton's
compactness theorem (Theorem 4.1.5) we see that the solution
$g_{ij}(x,t)$ of the modified flow actually converges
exponentially in $C^\infty$ topology to $g_{ij}(\infty)$. Moreover
the limiting metric $g_{ij}(\infty)$ satisfies
$$
M_{ij}=(r-R)g_{ij}+2\nabla_i\nabla_j\varphi=0,\ \ \mbox{ on }\; M.
$$
That is, the limiting metric is a shrinking gradient Ricci soliton
on the surface $M$.

The next result was first obtained by Hamilton in \cite{Ha88}. The
following simplified proof by using the Kazdan-Warner identity has
been widely known to experts in the field (e.g., cf. Proposition
5.21 of \cite{CK04}).

%\vskip 0.2cm \noindent {\bf Proposition 5.1.10} \emph{\
%CZ-SOURCE-LINE-10762
\begin{proposition} [rigidity of compact shrinking surface solitons]
  \label{prop:cz-compact-shrinking-surface-soliton-rigidity}
  \dcref{Cao--Zhu Proposition 5.1.10}
  \group{Normalized Ricci Flow on Surfaces}
  \source{cao-zhu}

On a compact surface there are no shrinking Ricci solitons other
than constant curvature.
\end{proposition}

%\vskip 0.1cm\noindent{\bf Proof.}\ \
\begin{proof}
By definition, a shrinking Ricci soliton on a compact surface $M$
is given by \begin{equation}
\nabla_iX_j+\nabla_jX_i=(R-r)g_{ij} %\eqno(5.1.26)
\end{equation} for some vector field $X={X_j}$. By contracting the above
equation by $Rg^{-1}$, we have
$$
2R(R-r)=2R\; \text{div}\, X,
$$
and hence
$$
\int_M(R-r)^2 dV=\int_MR(R-r)dV=\int_M R\; \text{div}\, X dV.
$$
Since $X$ is a conformal vector field (by the Ricci soliton
equation (5.1.26)), by integrating by parts and applying the
Kazdan-Warner identity \cite{KW}, we obtain
$$
\int_M(R-r)^2 dV=-\int_M\langle\nabla R,X\rangle dV=0.
$$
Hence $R\equiv r$, and the lemma is proved.
%$$\eqno \#$$ \vskip 0.3cm
\end{proof}

Now back to the solution of the modified flow (5.1.25). We have
seen the curvature converges exponentially to its limiting value
in the $C^\infty$ topology. But since there are no nontrivial
soliton on $M$, we must have $R$ converging exponentially to the
constant value $r$ in the $C^\infty$ topology. This then implies
that the unmodified flow (5.1.1) will converge to a metric of
positive constant curvature in the $C^\infty$ topology.

In conclusion, we have proved the following main theorem of this
section.

%\vskip 0.2cm\noindent{\bf Theorem 5.1.11} (
%CZ-SOURCE-LINE-10803
\begin{theorem}[convergence on surfaces of positive Euler characteristic]
  \label{thm:cz-positive-euler-surface-convergence}
  \dcref{Cao--Zhu Theorem 5.1.11}
  \group{Normalized Ricci Flow on Surfaces}
  \source{cao-zhu}

On a compact surface with $\chi(M)>0$, for any initial metric, the
solution of the normalized Ricci flow $(5.1.1)$ exists for all
time, and converges in the $C^\infty$ topology to a metric with
positive constant curvature.
%$$\eqno \#$$\vskip 1cm
\end{theorem}


\section{Differentiable Sphere Theorems in 3-D and 4-D}

An important problem in Riemannian geometry is to understand the
influence of curvatures, in particular the sign of curvatures, on
the topology of underlying manifolds. Classical results of this
type include sphere theorem and its refinements stated below
(e.g., cf. Theorem 6.1, Theorem 7.16, and Theorem 6.6 of
\cite{CE}). In this section we shall use the long-time behavior of
the Ricci flow on positively curved manifolds to establish
Hamilton's differentiable sphere theorems in dimensions three and
four. Our presentation is based on Hamilton \cite{Ha82, Ha86}.

Let us first recall some classical sphere theorems. Given a
Remannian manifold $M$, we denote by $K_M$ the sectional curvature
of $M$.

\medskip
{\bf Classical Sphere Theorems} (cf. \cite{CE}).\index{classical
sphere theorems} \ \emph{ Let $M$ be a complete, simply connected
$n$-dimensional manifold.}
\begin{itemize}
\item[(i)] \emph{If $\frac{1}{4}<K_M\le 1$, then $M$ is
homeomorphic to the $n$-sphere $\mathbb{S}^n$.} \item[(ii)]
\emph{There exists a positive constant $\delta\in(\frac{1}{4},1)$
such that if $\delta<K_M\le 1$, then $M$ is diffeomorphic to the
$n$-sphere $\mathbb{S}^n$.}
\end{itemize}

\medskip
Result (ii) is called the differentiable sphere theorem. If we
relax the assumptions on the strict lower bound in (i), then we
have the following rigidity result.

\medskip
{\bf Berger's Rigidity Theorem (cf. \cite{CE}).}\index{Berger's
rigidity theorem} \emph{ \ Let $M$ be a complete, simply connected
$n$-dimensional manifold with $\frac{1}{4}\le K_M\leq 1$. Then
either $M$ is homeomorphic to $\mathbb{S}^n$ or $M$ is isometric
to a symmetric space.}

\medskip
We remark that it follows from the classification of symmetric
spaces (see for example \cite{He}) that the only simply connected
symmetric spaces with positive curvature are $\mathbb{S}^n$,
$\mathbb{CP}^\frac{n}{2}$, $\mathbb{QP}^\frac{n}{4}$, and the
Cayley plane.

In early and mid 80's respectively, Hamilton \cite{Ha82, Ha86}
used the Ricci flow to prove the following differential sphere
theorems.

%\vskip 0.2cm \noindent {\bf Theorem 5.2.1} (
%CZ-SOURCE-LINE-10864
\begin{theorem}[three-dimensional positive-Ricci sphere theorem]
  \label{thm:cz-three-manifold-positive-ricci-sphere}
  \dcref{Cao--Zhu Theorem 5.2.1}
  \group{Differentiable Sphere Theorems}
  \source{cao-zhu}

A compact three-manifold with positive Ricci curvature must be
diffeomorphic to the three-sphere $\mathbb{S}^3$ or a quotient of
it by a finite group of fixed point free isometries in the
standard metric.
\end{theorem}

%\vskip 0.2cm \noindent {\bf Theorem 5.2.2} (
%CZ-SOURCE-LINE-10872
\begin{theorem}[four-dimensional positive-curvature-operator sphere theorem]
  \label{thm:cz-four-manifold-positive-curvature-operator-sphere}
  \dcref{Cao--Zhu Theorem 5.2.2}
  \group{Differentiable Sphere Theorems}
  \source{cao-zhu}

A compact four-manifold with positive curvature operator is
diffeomorphic to the four-sphere $\mathbb{S}^4$ or the real
projective space $\mathbb{RP}^4$.
\end{theorem}

Note that in above two theorems, we only assume curvatures to be
strictly positive, but not any strong pinching conditions as in
the classical sphere theorems. In fact, one of the important
special features discovered by Hamilton is that if the initial
metric has positive curvature, then the metric will get rounder
and rounder as it evolves under the Ricci flow, at least in
dimension three and four, so any small initial pinching will get
improved. Indeed, the pinching estimate is a key step in proving
both Theorem 5.2.1 and 5.2.2.

The following results are concerned with compact three-manifolds
or four-manifolds with weakly positive curvatures.

%\vskip 0.3cm \noindent {\bf Theorem 5.2.3} (Hamilton \cite{Ha86})
%CZ-SOURCE-LINE-10892
\begin{theorem}[classification under nonnegative Ricci or curvature operator]
  \label{thm:cz-nonnegative-curvature-low-dimensional-classification}
  \dcref{Cao--Zhu Theorem 5.2.3}
  \group{Differentiable Sphere Theorems}
  \source{cao-zhu}
 \
\begin{itemize}
\item[(i)] A compact three-manifold with nonnegative Ricci
curvature is diffeomorphic to $\mathbb{S}^3$, or a quotient of one
of the spaces $\mathbb{S}^3$ or $\mathbb{S}^2\times\mathbb{R}^1$
or $\mathbb{R}^3$ by a group of fixed point free isometries in the
standard metrics. \item[(ii)] A compact four-manifold with
nonnegative curvature operator is diffeomorphic to $\mathbb{S}^4$
or $\mathbb{CP}^2$ or $\mathbb{S}^2\times \mathbb{S}^2$, or a
quotient of one of the spaces $\mathbb{S}^4$ or $\mathbb{CP}^2$ or
$\mathbb{S}^3\times\mathbb{R}^1$ or $\mathbb{S}^2\times
\mathbb{S}^2$ or $\mathbb{S}^2\times \mathbb{R}^2$ or
$\mathbb{R}^4$ by a group of fixed point free isometries in the
standard metrics.
\end{itemize}
\end{theorem}

The rest of the section will be devoted to prove Theorems
5.2.1-5.2.3 and the presentation follows Hamilton \cite{Ha82,
Ha86} (also cf. \cite{Ha95F}).

Recall that the curvature operator $M_{\alpha\beta}$ evolves by
\begin{equation} \frac{\partial}{\partial t}M_{\alpha\beta}=\Delta
M_{\alpha\beta}+M^2_{\alpha\beta}+M_{\alpha\beta}^\#.
%\eqno(5.2.1)
\end{equation} where (see Section 1.3 and Section 2.4) $M_{\alpha\beta}^2$ is
the operator square
$$
M_{\alpha\beta}^2=M_{\alpha\gamma}M_{\beta\gamma}
$$
and $M_{\alpha\beta}^\#$ is the Lie algebra $so(n)$ square
$$
M_{\alpha\beta}^\#=C_\alpha^{\gamma\zeta}
C_\beta^{\eta\theta}M_{\gamma\eta}M_{\zeta\theta}.
$$

We begin with the curvature pinching estimates of the Ricci flow
in three dimensions. In dimension $n=3$, we know that
$M_{\alpha\beta}^\#$ is the adjoint matrix of $M_{\alpha\beta}$.
If we diagonalize $M_{\alpha\beta}$ with eigenvalues $\lambda \ge
\mu \ge \nu$ so that
$$
(M_{\alpha\beta})=\left(
\begin{array}{ccc}\lambda&  \ &  \ \\
\ &  \mu &  \ \\
\ &  \ &  \nu
\end{array}
\right),
$$
then $M_{\alpha\beta}^2$ and $M_{\alpha\beta}^\#$ are also
diagonal, with
$$
(M_{\alpha\beta}^2)=\left(
\begin{array}{ccc}\lambda^2&  \ &  \ \\
\ &  \mu^2 &  \ \\
\ &  \ &  \nu^2
\end{array}
\right) \ \ \text{and} \ \ (M_{\alpha\beta}^\#)=\left(
\begin{array}{ccc}\mu\nu&  \ &  \ \\
\ &  \lambda\nu &  \ \\
\ &  \ &  \lambda\mu
\end{array}
\right),
$$
and the ODE corresponding to PDE (5.2.1) is then given by the
system \begin{equation}
      \left\{
       \begin{array}{lll}
  \frac{d}{dt}\lambda=\lambda^2+\mu\nu,\\[4mm]
  \frac{d}{dt}\mu=\mu^2+\lambda\nu,\\[4mm]
  \frac{d}{dt}\nu=\nu^2+\lambda\mu.
       \end{array}
    \right.
 % \eqno(5.2.2)
\end{equation}

%\vskip 0.3cm \noindent {\bf Lemma 5.2.4}\ \ \emph{
%CZ-SOURCE-LINE-10969
\begin{lemma} [preservation of three-dimensional Ricci pinching]
  \label{lem:cz-three-dimensional-ricci-pinching-preserved}
  \dcref{Cao--Zhu Lemma 5.2.4}
  \group{Differentiable Sphere Theorems}
  \source{cao-zhu}

For any $\varepsilon\in[0,\frac{1}{3}]$, the pinching condition
$$
R_{ij}\geq0\ \ \ and\ \ \ R_{ij}\geq\varepsilon Rg_{ij}
$$
is preserved by the Ricci flow.
\end{lemma}

%\vskip 0.1cm \noindent {\bf Proof.}\ \
\begin{proof}
If we diagonalize the $3\times 3$ curvature operator matrix
$M_{\alpha\beta}$ with eigenvalues $\lambda \ge \mu \ge \nu$, then
nonnegative sectional curvature corresponds to $\nu\geq0$ and
nonnegative Ricci curvature corresponds to the inequality
$\mu+\nu\geq0$. Also, the scalar curvature $R=\lambda+\mu+\nu$. So
we need to show
$$
\mu+\nu\geq 0 \quad \text{and} \quad \mu+\nu\geq\delta \lambda, \;
\text{ with }\; \delta=2\varepsilon/(1-2\varepsilon),
$$
are preserved by the Ricci flow. By Hamilton's advanced maximum
principle (Theorem 2.3.1), it suffices to show that the closed
convex set
$$
K=\{M_{\alpha\beta}\ |\ \mu+\nu\geq0\; \text{ and }\;
\mu+\nu\geq\delta\lambda\}
$$
is preserved by the ODE system (5.2.2).

Now suppose we have diagonalized $M_{\alpha\beta}$ with
eigenvalues $\lambda \ge \mu \ge \nu$ at $t=0$, then both
$M^2_{\alpha\beta}$ and $M_{\alpha\beta}^\#$ are diagonal, so the
matrix $M_{\alpha\beta}$ remains diagonal for $t>0$. Moreover,
since
$$
\frac{d}{dt}(\mu-\nu)=(\mu-\nu)(\mu+\nu-\lambda),
$$
it is clear that $\mu\geq \nu$ for $t>0$ also. Similarly, we have
$\lambda\geq\mu$ for $t>0$. Hence the inequalities $\lambda \ge
\mu \ge \nu$ persist. This says that the solutions of the ODE
system (5.2.2) agree with the original choice for the eigenvalues
of the curvature operator.

The condition $\mu+\nu\geq 0$ is clearly preserved by the ODE,
because
$$
\frac{d}{dt}(\mu+\nu)=\mu^2+\nu^2+\lambda(\mu+\nu)\geq 0.
$$

It remains to check
$$
\frac{d}{dt}(\mu+\nu)\geq\delta\frac{d}{dt}\lambda
$$
or
$$
\mu^2+\lambda\nu+\nu^2+\lambda\mu\geq\delta(\lambda^2+\mu\nu)
$$
on the boundary where
$$
\mu+\nu=\delta \lambda\geq0.
$$
In fact, since
$$
(\lambda-\nu)\mu^2+(\lambda-\mu)\nu^2\geq0,
$$
we have
$$
\lambda(\mu^2+\nu^2)\geq(\mu+\nu)\mu\nu.
$$
Hence
\begin{displaymath}
\begin{split}
    \mu^2+\mu\lambda+\nu^2+\nu\lambda
    &  \geq\left(\frac{\mu+\nu}{\lambda}\right)(\lambda^2+\mu\nu)\\[3mm]
    &  =\delta(\lambda^2+\mu\nu).
\end{split}
\end{displaymath}
So we get the desired pinching estimate.
\end{proof}
%$$\eqno \#$$

%\vskip 0.3cm\noindent {\bf Proposition 5.2.5}\ \ \emph{
%CZ-SOURCE-LINE-11051
\begin{proposition} [asymptotic roundness under positive Ricci curvature]
  \label{prop:cz-three-dimensional-positive-ricci-asymptotic-roundness}
  \dcref{Cao--Zhu Proposition 5.2.5}
  \group{Differentiable Sphere Theorems}
  \source{cao-zhu}

Suppose that the initial metric of the solution to the Ricci flow
on $M^3\times [0,T)$ has positive Ricci curvature. Then for any
$\varepsilon>0$ we can find $C_\varepsilon<+\infty$ such that
$$
\left|R_{ij}-\frac{1}{3}Rg_{ij}\right|\leq\varepsilon
R+C_\varepsilon
$$
for all subsequent $t\in [0, T)$.
\end{proposition}

%\vskip 0.1cm \noindent {\bf Proof.}\ \ \
\begin{proof}
  \uses{lem:cz-three-dimensional-pinching-curvature-gradient}
  \uses{lem:cz-three-dimensional-ricci-pinching-preserved}
Again we consider the ODE system (5.2.2). Let $M_{\alpha\beta}$ be
diagonalized with eigenvalues $\lambda \ge \mu \ge \nu$ at $t=0$.
We saw in the proof of Lemma 5.2.4 the inequalities $\lambda \ge
\mu \ge \nu$ persist for $t>0$. We only need to show that there
are positive constants $\delta$ and $C$ such that the closed
convex set
$$
K=\{\ M_{\alpha\beta}\ |\ \lambda-\nu\leq
C(\lambda+\mu+\nu)^{1-\delta}\}
$$
is preserved by the ODE.

We compute
$$
\frac{d}{dt}(\lambda-\nu)=(\lambda-\nu)(\lambda+\nu-\mu)
$$
and
\begin{displaymath}
\begin{split}
   \frac{d}{dt}(\lambda+\mu+\nu)
& = (\lambda+\mu+\nu)(\lambda+\nu-\mu)+\mu^2\\[2mm]
&\quad  +\mu(\mu+\nu)+\lambda(\mu-\nu)\\[2mm]
& \geq (\lambda+\mu+\nu)(\lambda+\nu-\mu)+\mu^2.
\end{split}
\end{displaymath}
Thus, without loss of generality, we may assume $\lambda-\nu>0$
and get
$$
\frac{d}{dt}\log(\lambda-\nu)=\lambda+\nu-\mu
$$
and
$$
\frac{d}{dt}\log(\lambda+\mu+\nu)\geq
\lambda+\nu-\mu+\frac{\mu^2}{\lambda+\mu+\nu}.
$$
By Lemma 5.2.4, there exists a positive constant $C$ depending
only on the initial metric such that
$$
\lambda\leq \lambda+\mu\leq C(\mu+\nu)\leq 2C\mu,
$$
$$
\lambda+\nu-\mu\leq \lambda+\mu+\nu\leq 6C\mu,
$$
and hence with $\epsilon={1}/{36C^2}$,
$$
\frac{d}{dt}\log(\lambda+\mu+\nu)\geq(1+\epsilon)(\lambda+\nu-\mu).
$$
Therefore with $(1-\delta)={1}/{(1+\epsilon)}$,
$$
\frac{d}{dt}\log((\lambda-\nu)/(\lambda+\mu+\nu)^{1-\delta})\leq0.
$$
This proves the proposition.
\end{proof}
%$$\eqno\#$$ \vskip 0.3cm

We now are ready to prove Theorem 5.2.1.

\medskip
{\bf\em Proof of Theorem} {\bf 5.2.1.} \ Let $M$ be a compact
three-manifold with positive Ricci curvature and let the metric
evolve by the Ricci flow. By Lemma 5.2.4 we know that there exists
a positive constant $\beta>0$ such that
$$
R_{ij}\ge \beta Rg_{ij}
$$
for all $t\ge 0$ as long as the solution exists. The scalar
curvature evolves by
\begin{align*}
\frac{\partial R}{\partial t}&  = \Delta R+2|R_{ij}|^2\\
&  \ge \Delta R+\frac{2}{3}R^2,
\end{align*}
which implies, by the maximum principle, that the scalar curvature
remains positive and tends to $+\infty$ in finite time.

We now use a blow up argument as in Section 4.3 to get the
following gradient estimate.

\medskip
{\bf Claim.} \ For any $\varepsilon>0$, there exists a positive
constant $C_\varepsilon<+\infty$ such that for any time $\tau\ge
0$, we have
$$
\max_{t\le\tau}\max_{x\in M}|\nabla Rm(x,t)| \le \varepsilon
\max_{t\le\tau}\max_{x\in M}|Rm(x,t)|^\frac{3}{2}+C_\varepsilon.
$$

\medskip
We argue by contradiction. Suppose the above gradient estimate
fails for some fixed $\varepsilon_0>0$. Pick a sequence
$C_j\rightarrow+\infty$, and pick points $x_j\in M$ and times
$\tau_j$ such that
$$
|\nabla Rm(x_j,\tau_j)|\ge\varepsilon_0\max_{t\le\tau_j}
\max_{x\in M}|Rm(x,t)|^\frac{3}{2}+C_j,\quad j=1,2,\ldots.
$$
Choose $x_j$ to be the origin, and pull the metric back to a small
ball on the tangent space $T_{x_j}M$ of radius $r_j$ proportional
to the reciprocal of the square root of the maximum curvature up
to time $\tau_j$ (i.e., $\max_{t\le\tau_j}\max_{x\in M}|Rm(x,$
$t)|$). Clearly the maximum curvatures go to infinity by Shi's
derivative estimate of curvature (Theorem 1.4.1). Dilate the
metrics so that the maximum curvature
$$
\max_{t\le\tau_j}\max_{x\in M}|Rm(x,t)|
$$
becomes 1 and translate time so that $\tau_j$ becomes the time 0.
By Theorem 4.1.5, we can take a (local) limit. The limit metric
satisfies
$$
|\nabla Rm(0,0)|\ge \varepsilon_0>0.
$$
However the pinching estimate in Proposition 5.2.5 tells us the
limit metric has
$$
R_{ij}-\frac{1}{3}Rg_{ij}\equiv 0.
$$
By using the contracted second Bianchi identity
$$
\frac{1}{2}\nabla_iR =\nabla^jR_{ij}
=\nabla^j\left(R_{ij}-\frac{1}{3}Rg_{ij}\right)+\frac{1}{3}\nabla_iR,
$$
we get
$$
\nabla_iR\equiv0\quad \text{and then}\quad \nabla_iR_{jk}\equiv0.
$$
For a three-manifold, this in turn implies
$$
\nabla Rm=0
$$
which is a contradiction. Hence we have proved the gradient
estimate claimed.

We can now show that the solution to the Ricci flow becomes round
as the time $t$ tends to the maximal time $T$. We have seen that
the scalar curvature goes to infinity in finite time. Pick a
sequence of points $x_j\in M$ and times $\tau_j$ where the
curvature at $x_j$ is as large as it has been anywhere for $0\le
t\le\tau_j$ and $\tau_j$ tends to the maximal time. Since $|\nabla
Rm|$ is very small compared to $|Rm(x_j,\tau_j)|$ by the above
gradient estimate and $|R_{ij}-\frac{1}{3}Rg_{ij}|$ is also very
small compared to $|Rm(x_j,\tau_j)|$ by Proposition 5.2.5, the
curvature is nearly constant and positive in a large ball around
$x_j$ at the time $\tau_j$. But then the Bonnet-Myers' theorem
tells us this is the whole manifold. For $j$ large enough, the
sectional curvature of the solution at the time $\tau_j$ is
sufficiently pinched. Then it follows from the Klingenberg
injectivity radius estimate (see Section 4.2) that the injectivity
radius of the metric at time $\tau_j$ is bounded from below by
$c/\sqrt{|Rm(x_j,\tau_j)|}$ for some positive constant $c$
independent of $j$. Dilate the metrics so that the maximum
curvature $|Rm(x_j,\tau_j)|=\max_{t\leq\tau_j}\max_{x \in M}
|Rm(x,t)|$ becomes $1$ and shift the time $\tau_j$ to the new time
$0$. Then we can apply Hamilton's compactness theorem (Theorem
4.1.5) to take a limit. By the pinching estimate in Proposition
5.2.5, we know that the limit has positive constant curvature
which is either the round $\mathbb{S}^3$ or a metric quotient of
the round $\mathbb{S}^3$. Consequently, the compact three-manifold
$M$ is diffeomorphic to the round $\mathbb{S}^3$ or a metric
quotient of the round $\mathbb{S}^3$.
%$$\eqno \#$$
\endproof

Next we consider the pinching estimates of the Ricci flow on a
compact four-manifold $M$ with positive curvature operator. The
arguments are taken Hamilton \cite{Ha86}.

In dimension 4, we saw in Section 1.3 when we decompose
orthogonally $\Lambda^2=\Lambda^2_+\oplus\Lambda^2_-$ into the
eigenspaces of Hodge star with eigenvalue $\pm1$, we have a block
decomposition of $M_{\alpha\beta}$ as
$$
M_{\alpha\beta}=\left(\begin{array}{cc} A&  B\\ {}^tB&
C\end{array}\right)
$$
and then
$$
M_{\alpha\beta}^\#=2\left(\begin{array}{cc} A^\#&  B^\#\\
{}^tB^\#&  C^\#\end{array}\right)
$$
where $A^\#$, $B^\#$, $C^\#$ are the adjoints of $3\times3$
submatrices as before.

Thus the ODE
$$
\frac{d}{dt}M_{\alpha\beta}=M_{\alpha\beta}^2+M_{\alpha\beta}^\#
$$
corresponding to the PDE (5.2.1) breaks up into the system of
three equations \begin{equation}
   \left\{
   \begin{array}{lll}
    \frac {d }{d t}A=A^2+B {^tB}+2A^\#,
         \\[3mm]
    \frac {d }{d t}B=AB+BC+2B^\#,
         \\[3mm]
    \frac {d }{d t}C=C^2+ {^tB}B+2C^\#.
\end{array}
\right. %%\eqno (5.2.3)
\end{equation} As shown in Section 1.3, by the Bianchi identity, we know that
$\operatorname{tr}\, A=\operatorname{tr}\, C.$ For the symmetric matrices $A$ and $C$, we can
choose an orthonormal basis $x_1,x_2,x_3$ of $\Lambda^2_+$ such
that
$$
 A=
{ \left(
  \begin{array}{ccc}
    a_1&   0 &   0\\[4mm]
    0 &   a_2 &   0\\[4mm]
    0 &   0 &   a_3
\end{array}
\right),}
$$
and an orthonormal basis $z_1,z_2,z_3$ of $\Lambda^2_-$ such that
$$
 C=
{ \left(
  \begin{array}{ccc}
    c_1&   0 &   0\\[4mm]
    0 &   c_2 &   0\\[4mm]
    0 &   0 &   c_3
\end{array}
\right).}
$$
For matrix $B$, we can choose orthonormal basis
$y_1^+,y_2^+,y_3^+$ of $\Lambda^2_+$ and $y_1^-,y_2^-,$ $y_3^-$ of
$\Lambda^2_-$ such that
$$
 B=
{ \left(
  \begin{array}{ccc}
    b_1&   0 &   0\\[4mm]
    0 &   b_2 &   0\\[4mm]
    0 &   0 &   b_3
\end{array}
\right).}
$$
with $0\leq b_1\leq b_2\leq b_3$. We may also arrange the
eigenvalues of $A$ and $C$ as $a_1\leq a_2\leq a_3$ and $c_1\leq
c_2\leq c_3$. In view of the advanced maximum principle Theorem
2.3.1, we only need to establish the pinching estimates for the
ODE (5.2.3).

Note that
\begin{align*}
a_1&=\inf\{A(x,x)\ |\  x\in \Lambda_+^2 \ \ \text{and} \ \ |x|=1\},\\
a_3&=\sup\{A(x,x)\ |\  x\in \Lambda_+^2 \ \ \text{and} \ \ |x|=1\},\\
c_1&=\inf\{C(z,z)\ |\  z\in \Lambda_-^2 \ \ \text{and} \ \ |z|=1\},\\
c_3&=\sup\{C(z,z)\ |\  z\in \Lambda_-^2 \ \ \text{and} \ \
|z|=1\}.
\end{align*}
We can compute their derivatives by Lemma 2.3.3 as follows: \begin{equation}
   \left\{
   \begin{array}{lll}
    \frac {d }{d t}a_1\geq a_1^2+b_1^2+2a_2a_3,
         \\[3mm]
    \frac {d }{d t}a_3\leq a_3^2+b_3^2+2a_1a_2,
         \\[3mm]
    \frac {d }{d t}c_1\geq c_1^2+b_1^2+2c_2c_3,
         \\[3mm]
    \frac {d }{d t}c_3\leq c_3^2+b_3^2+2c_1c_2.
\end{array}
\right. %\eqno (5.2.4)
\end{equation} We shall make the pinching estimates by using the functions
$b_2+b_3$ and $a-2b+c$, where $a=a_1+a_2+a_3=c=c_1+c_2+c_3$ and
$b=b_1+b_2+b_3$. Since
\begin{align*}
b_2+b_3&  =B(y_2^+,y_2^-)+B(y_3^+,y_3^-) \\
&  = \sup\{B(y^+,y^-)+B(\tilde{y}^+,\tilde{y}^-)\, | \,
y^+,\tilde{y}^+\in \Lambda_+^2 \ \ \text{with} \ \
|y^+|=|\tilde{y}^+|=1,\\
&\qquad y^+ \bot \tilde{y}^+, \;\text{ and } \; \
y^-,\tilde{y}^-\in \Lambda_-^2 \;\text{with } \;
|y^-|=|\tilde{y}^-|=1,  y^- \bot \tilde{y}^- \},
\end{align*}
We compute by Lemma 2.3.3,
\begin{align}
\frac{d}{dt}(b_2+b_3)& \leq
\frac{d}{dt}B(y_2^+,y_2^-)+\frac{d}{dt}B(y_3^+,y_3^-)\\
&  = AB(y_2^+,y_2^-)+BC(y_2^+,y_2^-)+2B^\#(y_2^+,y_2^-) \nonumber\\
&\quad +AB(y_3^+,y_3^-)+BC(y_3^+,y_3^-)+2B^\#(y_3^+,y_3^-) \nonumber\\
&  = b_2A(y_2^+,y_2^+)+b_2C(y_2^-,y_2^-)+2b_1b_3 \nonumber\\
&\quad +b_3A(y_3^+,y_3^+)+b_3C(y_3^-,y_3^-)+2b_1b_2 \nonumber\\
&  \leq a_2b_2+a_3b_3+b_2c_2+b_3c_3+2b_1b_2+2b_1b_3,\nonumber
\end{align} %\eqno(5.2.5)
where we used the facts that $A(y_2^+,y_2^+)+A(y_3^+,y_3^+)\leq
a_2+a_3$ and $C(y_2^-,y_2^-)+C(y_3^-,y_3^-)\leq c_2+c_3$.

Note also that the function $a=trA=c=trC$ is linear, and the
function $b$ is given by
\begin{align*}
b&  = B(y_1^+,y_1^-)+B(y_2^+,y_2^-)+B(y_3^+,y_3^-) \\
& =\sup\Big\{B(Ty_1^+,\tilde{T}y_1^-)+B(Ty_2^+,\tilde{T}y_2^-)
+B(Ty_3^+,\tilde{T}y_3^-)\ |
\ T,\tilde{T} \; \text{ are} \\
&\qquad \text{othogonal transformations of}\ \Lambda^2_+ \; \text{
and }\; \Lambda^2_-\; \text{ respectively}\Big\}.
\end{align*}
Indeed,
\begin{align*}
& B(Ty_1^+,\tilde{T}y_1^-)+B(Ty_2^+,\tilde{T}y_2^-)
+B(Ty_3^+,\tilde{T}y_3^-)\\
&  = B(y_1^+,T^{-1}\tilde{T}(y_1^-))
+B(y_2^+,T^{-1}\tilde{T}(y_2^-))+B(y_3^+,T^{-1}\tilde{T}(y_3^-))\\
&  = b_1t_{11}+b_2t_{22}+b_3t_{33}
\end{align*}
where $t_{11},t_{22},t_{33}$ are diagonal elements of the
orthogonal matrix $T^{-1}\tilde{T}$ with $t_{11},t_{22},t_{33}\leq
1$. Thus by using Lemma 2.3.3 again, we compute
\begin{align*}
\frac{d}{dt}(a-2b+c)&  \geq
{\rm tr}\,\left(\frac{d}{dt}A-2\frac{d}{dt}B+\frac{d}{dt}C\right)\\
&  ={\rm tr}\,((A-B)^2+(C-B)^2+2(A^\#-2B^\#+C^\#))
\end{align*}
evaluated in those coordinates where $B$ is diagonal as above.
Recalling the definition of Lie algebra product
$$
P\#Q=\frac{1}{2}\varepsilon_{\alpha\beta\gamma}
\varepsilon_{\zeta\eta\theta}P_{\beta\eta}Q_{\gamma\theta}
$$
with $\varepsilon_{\alpha\beta\gamma}$ being the permutation
tensor, we see that the Lie algebra product $\#$ gives a symmetric
bilinear operation on matrices, and then
\begin{align*}
& {\rm tr}\,(2(A^\#-2B^\#+C^\#)) \\
&  = {\rm tr}\,((A-C)^\#+(A+2B+C)\#(A-2B+C))\\
&  = -\frac{1}{2}{\rm tr}\,(A-C)^2+\frac{1}{2}({\rm tr}\,(A-C))^2\\
&\quad + {\rm tr}\,((A+2B+C)\#(A-2B+C))\\
&  = -\frac{1}{2}{\rm tr}\,(A-C)^2+{\rm tr}\,((A+2B+C)\#(A-2B+C))
\end{align*}
by the Bianchi identity. It is easy to check that
$$
{\rm tr}\,(A-B)^2+\operatorname{tr}\,(C-B)^2-\frac{1}{2}\operatorname{tr}\,(A-C)^2
=\frac{1}{2}\operatorname{tr}\,(A-2B+C)^2\geq 0.
$$
Thus we obtain
$$
\frac{d}{dt}(a-2b+c)\geq \operatorname{tr}\,((A+2B+C)\#(A-2B+C))
$$
Since $M_{\alpha\beta}\geq 0$ and
$$
M_{\alpha\beta}= { \left(
  \begin{array}{cc}
    A &   B\\[4mm]
    ^tB &   C
   \end{array}
\right),}
$$
we see that $A+2B+C\geq 0$ and $A-2B+C\geq 0$, by applying
$M_{\alpha\beta}$ to the vectors $(x,x)$ and $(x,-x)$. It is then
not hard to see
$$
\operatorname{tr}\,((A+2B+C)\#(A-2B+C))\geq (a_1+2b_1+c_1)(a-2b+c).
$$
Hence we obtain \begin{equation}
\frac{d}{dt}(a-2b+c)\geq (a_1+2b_1+c_1)(a-2b+c). %\eqno (5.2.6)
\end{equation}

We now state and prove the following pinching estimates of
Hamilton for the associated ODE (5.2.3).

%\vskip 0.2cm \noindent {\bf Proposition 5.2.6.} (
%CZ-SOURCE-LINE-11424
\begin{proposition}[invariant four-dimensional curvature pinching set]
  \label{prop:cz-four-dimensional-invariant-curvature-pinching-set}
  \dcref{Cao--Zhu Proposition 5.2.6}
  \group{Differentiable Sphere Theorems}
  \source{cao-zhu}

If we choose successively positive constants $G$ large enough, $H$
large enough, $\delta$ small enough, $J$ large enough,
$\varepsilon$ small enough, $K$ large enough, $\theta$ small
enough, and $L$ large enough, with each depending on those chosen
before, then the closed convex subset $X$ of
$\{M_{\alpha\beta}\geq 0 \}$ defined by the inequalities
\begin{itemize}
\item[(1)] $(b_2+b_3)^2\leq Ga_1c_1,$ \item[(2)] $a_3\leq Ha_1 \ \
\text{and}\ \ c_3\leq Hc_1,$ \item[(3)] $(b_2+b_3)^{2+\delta}\leq
Ja_1c_1(a-2b+c)^{\delta},$ \item[(4)]
$(b_2+b_3)^{2+\varepsilon}\leq Ka_1c_1,$ \item[(5)] $a_3\leq
a_1+La_1^{1-\theta}\ \ \text{and} \ \ c_3\leq
c_1+Lc_1^{1-\theta},$
\end{itemize}
is preserved by {\rm ODE} $(5.2.3)$. Moreover every compact subset
of $\{M_{\alpha\beta}>0\}$ lies in some such set $X$.
\end{proposition}

%\vskip 0.1cm \noindent {\bf Proof.} \
\begin{proof}
Clearly the subset $X$ is closed and convex.  We first note that
we may assume $b_2+b_3>0$ because if $b_2+b_3=0$, then from
(5.2.5), $b_2+b_3$ will remain zero and then the inequalities (1),
(3) and (4) concerning $b_2+b_3$ are automatically satisfied.
Likewise we may assume $a_3>0$ and $c_3>0$ from (5.2.4).

Let $G$ be a fixed positive constant. To prove the inequality (1)
we only need to check \begin{equation}
\frac{d}{dt}\log\frac{a_1c_1}{(b_2+b_3)^2}\geq 0 %\eqno (5.2.7)
\end{equation} whenever $(b_2+b_3)^2=Ga_1c_1$ and $b_2+b_3>0$. Indeed, it
follows from (5.2.4) and (5.2.5) that
\begin{align}
\frac{d}{dt}\log a_1 &\geq
2b_1+2a_3+\frac{(a_1-b_1)^2}{a_1}+2\frac{a_3}{a_1}(a_2-a_1), \\
%\eqno(5.2.8)
\frac{d}{dt}\log c_1 &\geq
2b_1+2c_3+\frac{(c_1-b_1)^2}{c_1}+2\frac{c_3}{c_1}(c_2-c_1), \\
%\eqno(5.2.9)
\intertext{and} \frac{d}{dt}\log(b_2+b_3)&\leq
2b_1+a_3+c_3-\frac{b_2}{b_2+b_3}[(a_3-a_2)+(c_3-c_2)],
%\eqno(5.2.10)
\end{align}
which immediately give the desired inequality (5.2.7).

By (5.2.4), we have \begin{equation} \frac{d}{dt}\log a_3\leq
a_3+2a_1+\frac{b_3^2}{a_3}-\frac{2a_1}{a_3}(a_3-a_2).
%\eqno(5.2.11)
\end{equation} {}From the inequality (1) there holds $b_3^2\leq Ga_1c_1$.
Since $\operatorname{tr}\, A=\operatorname{tr}\, C$, $c_1\leq c_1+c_2+c_3=a_1+a_2+a_3\leq 3a_3$
\noindent which shows
$$
\frac{b_3^2}{a_3}\leq 3Ga_1.
$$
Thus by (5.2.8) and (5.2.11),
$$
\frac{d}{dt}\log\frac{a_3}{a_1}\leq (3G+2)a_1-a_3.
$$
So if $H\geq (3G+2)$, then the inequalities $a_3\leq Ha_1$ and
likewise $c_3\leq Hc_1$ are preserved.

For the inequality (3), we compute from (5.2.8)-(5.2.10)
\begin{align*}
\frac{d}{dt}\log\frac{a_1c_1}{(b_2+b_3)^2}&  \geq
\frac{(a_1-b_1)^2}{a_1}+\frac{(c_1-b_1)^2}{c_1}
+2\frac{a_3}{a_1}(a_2-a_1)+2\frac{c_3}{c_1}(c_2-c_1)\\
&\quad +\frac{2b_2}{b_2+b_3}[(a_3-a_2)+(c_3-c_2)].
\end{align*}
If $b_1\leq a_1/2$, then
$$
\frac{(a_1-b_1)^2}{a_1}\geq \frac{a_1}{4}\geq\frac{1}{4H}a_3,
$$
and if $b_1\geq a_1/2$, then
$$
\frac{2b_2}{b_2+b_3}\geq \frac{2b_2}{\sqrt{Ga_1c_1}}\geq
\frac{2b_2}{\sqrt{3Ga_1a_3}}\geq \frac{2b_2}{\sqrt{3GH}\cdot a_1}
\geq \frac{1}{\sqrt{3GH}}.
$$
Thus by combining with $ 3a_3 \geq c_3$, we have
$$
\frac{d}{dt}\log\frac{a_1c_1}{(b_2+b_3)^2}\geq \delta
(a_3-a_1)+\delta(c_3-c_1)
$$
provided $\delta\leq \min(\frac{1}{24H},\frac{1}{\sqrt{3GH}})$. On
the other hand, it follows from (5.2.6) and (5.2.10) that
$$
\frac{d}{dt}\log\frac{b_2+b_3}{a-2b+c}\leq(a_3-a_1)+(c_3-c_1).
$$
Therefore the inequality (3)
$$
(b_2+b_3)^{2+\delta}\le Ja_1c_1(a-2b+c)^\delta
$$
will be preserved by any positive constant $J$.

To verify the inequality (4), we first note that there is a small
$\eta>0$ such that
$$
b\le (1-\eta)a,
$$
on the set defined by the inequality (3). Indeed, if $b\le
\frac{a}{2}$, this is trivial and if $b\ge\frac{a}{2}$, then
$$
\left(\frac{a}{3}\right)^{2+\delta} \le(b_2+b_3)^{2+\delta}\le 2^\delta
Ja^2(a-b)^\delta
$$
which makes $b\le(1-\eta)a$ for some $\eta>0$ small enough.
Consequently, $$\eta a\le a-b\le 3(a_3-b_1)$$ which implies either
$$
a_3-a_1\ge \frac{1}{6}\eta a,
$$
or
$$
a_1-b_1\ge \frac{1}{6}\eta a.
$$
Thus as in the proof of the inequality (3), we have
$$
\frac{d}{dt}\log\frac{a_1c_1}{(b_2+b_3)^2}\ge \delta(a_3-a_1)
$$
and
$$
\frac{d}{dt}\log\frac{a_1c_1}{(b_2+b_3)^2}\ge
\frac{(a_1-b_1)^2}{a_1},
$$
which in turn implies
$$
\frac{d}{dt}\log\frac{a_1c_1}{(b_2+b_3)^2} \ge
\left(\max\left\{\frac{1}{6}\eta\delta,\;
\frac{1}{36}\eta^2\right\}\right)\cdot a.
$$
On the other hand, it follows from (5.2.10) that
$$
\frac{d}{dt}\log(b_2+b_3)\le 2b_1+a_3+c_3\le 4a
$$
since $M_{\alpha\beta}\ge 0$. Then if $\varepsilon>0$ is small
enough
$$
\frac{d}{dt}\log\frac{a_1c_1}{(b_2+b_3)^{2+\varepsilon}}\ge 0
$$
and it follows that the inequality (4) is preserved by any
positive $K$.

Finally we consider the inequality (5). From (5.2.8) we have
$$
\frac{d}{dt}\log a_1\ge a_1+2a_3
$$
and then for $\theta\in (0,1)$,
$$
\frac{d}{dt}\log (a_1+La_1^{1-\theta})\ge
\frac{a_1+(1-\theta)La_1^{1-\theta}}{a_1+La_1^{1-\theta}}(a_1+2a_3).
$$
On the other hand, the inequality (4) tells us
$$
b_3^2\le \tilde{K}a_1^{1-\theta}a_3
$$
for some positive constant $\tilde{K}$ large enough with $\theta$
to be fixed small enough. And then
$$
\frac{d}{dt}\log a_3\le a_3+2a_1+\tilde{K}a_1^{1-\theta},
$$
by combining with (5.2.11). Thus by choosing
$\theta\le\frac{1}{6H}$ and $L\ge 2\tilde{K}$,
\begin{align*}
\frac{d}{dt}\log \frac{a_1+La_1^{1-\theta}}{a_3}& \ge
(a_3-a_1)-\theta\frac{La_1^{1-\theta}}{a_1+La_1^{1-\theta}}(a_1+2a_3)
-\tilde{K}a_1^{1-\theta}\\
&
\ge(a_3-a_1)-\theta\frac{La_1^{1-\theta}}{a_1+La_1^{1-\theta}}\cdot
3Ha_1-\tilde{K}a_1^{1-\theta}\\
&  \ge (a_3-a_1)-(3\theta HL+\tilde{K})a_1^{1-\theta}\\
&  = [L-(3\theta HL+\tilde{K})]a_1^{1-\theta}\\
&  \ge 0
\end{align*}
whenever $a_1+La_1^{1-\theta}=a_3$. Consequently the set
$\{a_1+La_1^{1-\theta}\ge a_3\}$ is preserved. A similar argument
works for the inequality in $C$. This completes the proof of
Proposition 5.2.6.
%$$\eqno \#$$ \vskip 0.3cm
\end{proof}

The combination of the advanced maximum principle Theorem 2.3.1
and the pinching estimates of the ODE (5.2.3) in Proposition 5.2.6
immediately gives the following pinching estimate for the Ricci
flow on a compact four-manifold.

%\vskip 0.2cm \noindent{\bf Corollary 5.2.7} $\;\;${\sl
%CZ-SOURCE-LINE-11609
\begin{corollary} [four-dimensional curvature-operator asymptotic roundness]
  \label{cor:cz-four-dimensional-curvature-operator-roundness}
  \dcref{Cao--Zhu Corollary 5.2.7}
  \group{Differentiable Sphere Theorems}
  \source{cao-zhu}

Suppose that the initial metric of the solution to the Ricci flow
on a compact four-manifold has positive curvature operator. Then
for any $\varepsilon>0$ we can find positive constant
$C_\varepsilon<+\infty$ such that
$$
|{\mathop{Rm}^ \circ}|\le\varepsilon R+C_\varepsilon
$$
for all $t\ge 0$ as long as the solution exists, where ${\mathop
{Rm}\limits^ \circ}$ is the traceless part of the curvature
operator.
%$$\eqno \#$$ \vskip 0.3cm
\end{corollary}

\smallskip
{\bf\em Proof of Theorem} {\bf 5.2.2.} \ Let $M$ be a compact
four-manifold with positive curvature operator and let us evolve
the metric by the Ricci flow. Again the evolution equation of the
scalar curvature tells us that the scalar curvature remains
positive and becomes unbounded in finite time. Pick a sequence of
points $x_j\in M$ and times $\tau_j$ where the curvature at $x_j$
is as large as it has been anywhere for $0\le t\le \tau_j$. Dilate
the metrics so that the maximum curvature
$|Rm(x_j,\tau_j)|=\max_{t\le\tau_j}\max_{x\in M}|Rm(x,t)|$ becomes
1 and shift the time so that the time $\tau_j$ becomes the new
time 0. The Klingenberg injectivity radius estimate in Section 4.2
tells us that the injectivity radii of the rescaled metrics at the
origins $x_j$ and at the new time 0 are uniformly bounded from
below. Then we can apply the Hamilton's compactness theorem
(Theorem 4.1.5) to take a limit. By the pinching estimate in
Corollary 5.2.7, we know that the limit metric has positive
constant curvature which is either $\mathbb{S}^4$ or
$\mathbb{R}\mathbb{P}^4$. Therefore the compact four-manifold $M$
is diffeomorphic to the sphere $\mathbb{S}^4$ or the real
projective space $\mathbb{R}\mathbb{P}^4$.
%$$\eqno \#$$ \vskip 0.3cm
\endproof

%\vskip 0.3cm \noindent \textbf{Remark 5.2.8.}
%CZ-SOURCE-LINE-11648
\begin{remark}[subsequential convergence in the differentiable sphere theorems]
  \label{rem:cz-differentiable-sphere-subsequential-convergence}
  \dcref{Cao--Zhu Remark 5.2.8}
  \group{Differentiable Sphere Theorems}
  \source{cao-zhu}
  \uses{thm:cz-three-manifold-positive-ricci-sphere, thm:cz-four-manifold-positive-curvature-operator-sphere}

The proofs of Theorem 5.2.1 and Theorem 5.2.2 also show that the
Ricci flow on a compact three-manifold with positive Ricci
curvature or a compact four-manifold with positive curvature
operator is subsequentially converging (up to scalings) in the
$C^{\infty}$ topology to the same underlying compact manifold with
a metric of positive constant curvature. Of course, this
subsequential convergence is in the sense of Hamilton's
compactness theorem (Theorem 4.1.5) which is also up to the
pullbacks of diffeomorphisms. Actually in \cite{Ha82} and
\cite{Ha86}, Hamilton obtained the convergence in the stronger
sense that the (rescaled) metrics converge (in the $C^{\infty}$
topology) to a constant (positive) curvature metric.
%\vskip 0.3cm
\end{remark}

In the following we use Hamilton's strong maximum principle
(Theorem 2.2.1) to prove Theorem 5.2.3.

\medskip
{\bf \em Proof of Theorem} {\bf 5.2.3.} \ In views of Theorem
5.2.1 and Theorem 5.2.2, we may assume the Ricci curvature (in
dimension 3) and the curvature operator (in dimension 4) always
have nontrivial kernels somewhere along the Ricci flow.

\smallskip
(i) In the case of dimension 3, we consider the evolution equation
(1.3.5) of the Ricci curvature
$$
\frac{\partial R_{ab}}{\partial t} =\triangle
R_{ab}+2R_{acbd}R_{cd}
$$
in an orthonormal frame coordinate. At each point, we diagonalize
$R_{ab}$ with eigenvectors $e_{1}, e_{2}, e_{3}$ and eigenvalues
$\lambda_1\leq\lambda_2\leq\lambda_3$. Since
\begin{align*}
R_{1c1d}R_{cd}&  = R_{1212}R_{22}+R_{1313}R_{33}\\
&  =\frac{1}{2}((\lambda_{3}-\lambda_{2})^{2}
+\lambda_{1}(\lambda_2+\lambda_3)),
\end{align*}
we know that if $R_{ab}\geq 0$, then $R_{acbd}R_{cd}\geq0$. By
Hamilton's strong maximum principle (Theorem 2.2.1), there exists
an interval $0<t<\delta$ on which the rank of $R_{ab}$ is constant
and the null space of $R_{ab}$ is invariant under parallel
translation and invariant in time and also lies in the null space
of $R_{acbd}R_{cd}$. If the null space of $R_{ab}$ has rank one,
then $\lambda_1=0$ and $\lambda_2=\lambda_{3}>0$.  In this case,
by De Rham decomposition theorem, the universal cover $\tilde{M}$
of the compact $M$ splits isometrically as $\mathbb{R}\times
\Sigma^{2}$ and the curvature of $\Sigma^{2}$ has a positive lower
bound. Hence $\Sigma^{2}$ is diffeomorphic to $\mathbb{S}^{2}$.
Assume $M=\mathbb{R}\times \Sigma^{2}/\Gamma$, for some isometric
subgroup $\Gamma$ of $\mathbb{R}\times \Sigma^{2}$. Note that
$\Gamma$ remains to be an isometric subgroup of $\mathbb{R}\times
\Sigma^{2}$ during the Ricci flow by the uniqueness (Theorem
1.2.4). Since the Ricci flow on $\mathbb{R}\times
\Sigma^{2}/\Gamma$ converges to the standard metric by Theorem
5.1.11, $\Gamma$ must be an isometric subgroup of
$\mathbb{R}\times\mathbb{S}^{2}$ in the standard metric. If the
null space of $R_{ab}$ has rank greater than one, then $R_{ab}=0$
and the manifold is flat. This proves Theorem 5.2.3 part (i).

\smallskip
(ii) In the case of dimension 4, we classify the manifolds
according to the (restricted) holonomy algebra $\mathcal{G}$. Note
that the curvature operator has nontrivial kernel and
$\mathcal{G}$ is the image of the the curvature operator, we see
that $\mathcal{G}$ is a proper subalgebra of $so(4)$. We divide
the argument into two cases.

\smallskip
{\it Case} 1. $\mathcal{G}$ is reducible.

\smallskip
In this case the universal cover $\tilde{M}$ splits isometrically
as $\tilde{M_{1}}\times \tilde{M_{2}}$.  By the above results on
two and three dimensional Ricci flow, we see that ${M}$ is
diffeomorphic to a quotient of one of the spaces $\mathbb{R}^{4},$
$\mathbb{R}\times \mathbb{S}^{3}$, $\mathbb{R}^{2}\times
\mathbb{S}^{2}$, $\mathbb{S}^{2}\times \mathbb{S}^{2}$ by a group
of fixed point free isometries. As before by running the Ricci
flow until it converges and using the uniqueness (Theorem 1.2.4),
we see that this group is actually a subgroup of the isometries in
the standard metrics.

\smallskip
{\it Case} 2. $\mathcal{G}$ is not reducible (i.e., irreducible).

\smallskip
If the manifold is not Einstein, then by Berger's classification
theorem for holonomy groups \cite{Ber}, $\mathcal{G}=so(4)$ or
$u(2)$. Since the curvature operator is not strictly positive,
$\mathcal{G}=u(2)$, and the universal cover $\tilde{M}$ of $M$ is
K$\ddot{a}$hler and has positive bisectional curvature. In this
case $\tilde{M}$ is biholomorphic to $\mathbb{C}\mathbb{P}^{2}$ by
the result of Andreotti-Frankel \cite{Ft} (also cf. Mori
\cite{Mori} and Siu-Yau \cite{SiY}).

If the manifold is Einstein, then by the block decomposition of
the curvature operator matrix in four-manifolds (see the third
section of Chapter 1),
$$
Rm(\Lambda^{2}_{+},\Lambda^{2}_{-})=0.
$$
Let $\varphi\neq0$, and
$$
\varphi=\varphi_{+}+\varphi_{-}\in
  \Lambda^{2}_{+}\oplus\Lambda^{2}_{-},
$$
lies in the kernel of the curvature operator, then
$$
   0=Rm(\varphi_{+},\varphi_{+})+Rm(\varphi_{-},\varphi_{-}).
$$
It follows that
$$
   Rm(\varphi_{+},\varphi_{+})=0, \mbox{ and }
   Rm(\varphi_{-},\varphi_{-})=0.
$$
We may assume $\varphi_{+}\neq 0$ (the argument for the other case
is similar).  We consider the restriction of $Rm$ to
$\Lambda^{2}_{+}$, since $\Lambda^{2}_{+}$ is an invariant
subspace of $R{m}$ and the intersection of $\Lambda^{2}_{+}$ with
the null space of $Rm$ is nontrivial. By considering the null
space of $Rm$ and its orthogonal complement in $\Lambda^{2}_{+}$,
we obtain a parallel distribution of rank one in
$\Lambda^{2}_{+}$. This parallel distribution gives a parallel
translation invariant two-form $\omega \in \Lambda^{2}_{+}$ on the
universal cover $\tilde{M}$ of $M$. This two-form is
nondegenerate, so it induces a K$\ddot{a}$hler structure of
$\tilde{M}$. Since the K$\ddot{a}$hler metric is parallel with
respect to the original metric and the manifold is irreducible,
the K$\ddot{a}$hler metric is proportional to the original metric.
Hence the manifold $\tilde{M}$ is K$\ddot{a}$hler-Einstein with
nonnegative curvature operator. Taking into account the
irreducibility of $\mathcal{G}$, it follows that $\tilde{M}$ is
biholomorphic to $\mathbb{C}\mathbb{P}^{2}$.  Therefore the proof
of Theorem 5.2.3 is completed.
%$$\eqno \#$$ \vskip 0.3cm
\endproof

To end this section, we mention some generalizations of Hamilton's
differential sphere theorem (Theorem 5.2.1 and Theorem 5.2.2) to
higher dimensions.


It is well-known that the curvature tensor $Rm=\{R_{ijkl}\}$ of a
Riemannian manifold can be decomposed into three orthogonal
components which have the same symmetries as $Rm$:
$$
Rm=W+V+U.
$$
Here $W=\{W_{ijkl}\}$ is the Weyl conformal curvature tensor,
whereas $V=\{V_{ijkl}\}$ and $U=\{U_{ijkl}\}$ denote the traceless
Ricci part and the scalar curvature part respectively. The
following pointwisely pinching sphere theorem under the additional
assumption that the manifold is compact was first obtained by
Huisken \cite{Hu}, Margerin \cite{Mar94}, \cite{Mar98} and
Nishikawa \cite{Ns} by using the Ricci flow. The compactness
assumption was later removed by Chen and the second author in
\cite{CZ00}.

% \vskip 0.3cm \noindent {\bf Theorem 5.2.9} \
%CZ-SOURCE-LINE-11810
\begin{theorem}[higher-dimensional pointwise-pinched sphere theorem]
  \label{thm:cz-higher-dimensional-pointwise-pinched-sphere}
  \dcref{Cao--Zhu Theorem 5.2.9}
  \group{Differentiable Sphere Theorems}
  \source{cao-zhu}

Let $n\ge4$. Suppose $M$ is a complete n-dimensional manifold with
positive and bounded scalar curvature and satisfies the
pointwisely pinching condition
$$
|W|^2+|V|^2\le\delta_n(1-\varepsilon)^2|U|^2,
$$
where $\varepsilon>0,\delta_4=\frac{1}{5},\delta_5=\frac{1}{10}$,
and
$$
\delta_n=\frac{2}{(n-2)(n+1)},n\ge 6.
$$
Then $M$ is diffeomorphic to the sphere $\mathbb{S}^n$ or a
quotient of it by a finite group of fixed point free isometries in
the standard metric.
%$$\eqno \#$$ \vskip 0.3cm
\end{theorem}

Also, using the minimal surface theory, Micallef and Moore
\cite{MiMo} proved any compact simply connected $n$-dimensional
manifold with positive curvature operator is
homeomorphic\footnote{Very recently, B$\ddot{o}$hm and Wilking
\cite{BW} have proved, by using the Ricci flow, that a compact
simply connected $n$-dimensional manifold with positive (or
2-positive) curvature operator is {\it diffeomorphic} to
$\mathbb{S}^n$. This gives an affirmative answer to a
long-standing conjecture of Hamilton.} to the $n$-sphere
$\mathbb{S}^n$.

Finally, in \cite{CZ00}, Chen and the second author also used the
Ricci flow to obtain the following flatness theorem for noncompact
three-manifolds.

%\vskip 0.3cm\noindent {\bf Theorem 5.2.10} \ \emph{
%CZ-SOURCE-LINE-11844
\begin{theorem}[flatness from Ricci pinching in the noncompact three-dimensional case]
  \label{thm:cz-noncompact-three-manifold-ricci-pinching-flatness}
  \dcref{Cao--Zhu Theorem 5.2.10}
  \group{Differentiable Sphere Theorems}
  \source{cao-zhu}

Let $M$ be a three-dimensional complete noncompact Riemannian
manifold with bounded and nonnegative sectional curvature. Suppose
$M$ satisfies the following Ricci pinching condition
$$
R_{ij}\ge \varepsilon Rg_{ij},\quad \text{on }\; M,
$$
for some $\varepsilon>0$. Then $M$ is flat.
%$$\eqno \#$$ \vskip 0.3cm
\end{theorem}

\vskip 1cm


\section{Nonsingular Solutions on Three-manifolds}

We have seen in the previous section that a good understanding of
the long time behaviors for solutions to the Ricci flow could lead
to remarkable topological or geometric consequences for the
underlying manifolds. Since one of the central themes of the Ricci
flow is to study the geometry and topology of three-manifolds, we
will start to analyze the long time behavior of the Ricci flow on
a compact three-manifold. Here, we shall first consider a special
class of solutions, the nonsingular solutions (see the definition
below). The main purpose of this section is to present Hamilton's
important result in \cite{Ha99} that any compact three-manifold
admitting a nonsingular solution is geometrizable in the sense of
Thurston\cite{Th82}. Most of the presentation is based on Hamilton
\cite{Ha99}.

Let $M$ be a compact three-manifold. We will consider the
(unnormalized) Ricci flow
$$
\frac{\partial}{\partial t}g_{ij}=-2R_{ij},
$$
and the normalized Ricci flow
$$
\frac{\partial}{\partial t}g_{ij}=\frac{2}{3}rg_{ij}-2R_{ij}
$$
where $r=r(t)$ is the function of the average of the scalar
curvature. Recall that the normalized flow differs from the
unnormalized flow only by rescaling in space and time so that the
total volume $V=\int_Md\mu$ remains constant. As we mentioned
before, in this section we only consider a special class of
solutions that we now define.

%\vskip 0.2cm \noindent {\bf Definition 5.3.1}\ \
%CZ-SOURCE-LINE-11891
\begin{definition}[nonsingular normalized Ricci-flow solution]
  \label{def:cz-nonsingular-normalized-ricci-flow}
  \dcref{Cao--Zhu Definition 5.3.1}
  \group{Nonsingular Three-Dimensional Solutions}
  \source{cao-zhu}

A {\bf nonsingular solution}\index{solution!nonsingular} of the
Ricci flow is one where the solution of the normalized flow exists
for all time $0\leq t<\infty$, and the curvature remains bounded
$|Rm|\leq C<+\infty$ for all time with some constant $C$
independent of $t$.
\end{definition}

Clearly any solution to the Ricci flow on a compact three-manifold
with nonnegative Ricci curvature is nonsingular. Currently there
are few conditions which guarantee a solution will remain
nonsingular. Nevertheless, the ideas and arguments of Hamilton
\cite{Ha99} as described below is extremely important. One will
see in Chapter 7 that these arguments will be modified to analyze
the long-time behavior of arbitrary solutions, or even the
solutions with surgery, to the Ricci flow on three-manifolds.

We begin with an improvement of Hamilton-Ivey pinching result
\index{Hamilton-Ivey pinching estimate} (Theorem 2.4.1).

%\vskip 0.2cm \noindent {\bf Theorem 5.3.2}\
%CZ-SOURCE-LINE-11912
\begin{theorem}[Hamilton--Ivey pinching estimate]
  \label{thm:cz-hamilton-ivey-pinching-estimate}
  \dcref{Cao--Zhu Theorem 5.3.2}
  \group{Nonsingular Three-Dimensional Solutions}
  \source{cao-zhu}

Suppose we have a solution to the $($unnormalized$\,)$\; Ricci\;
flow\; on\; a\; three--manifold\; which\; is\; complete\; with
bounded curvature for each $t\geq0$. Assume at $t=0$ the
eigenvalues $\lambda\geq\mu\geq\nu$ of the curvature operator at
each point are bounded below by $\nu\geq-1$. Then at all points
and all times $t\geq0$ we have the pinching estimate
$$
R\geq(-\nu)[\log(-\nu)+\log(1+t)-3]
$$
whenever $\nu<0$.
\end{theorem}

%\vskip 0.1cm \noindent {\bf Proof.} \ \
\begin{proof}
As before, we study the ODE system
\begin{equation*}
\begin{cases}{\displaystyle
    \frac{d\lambda}{dt} =\lambda^2+\mu\nu,}\\[3mm]
{\displaystyle \frac{d\mu}{dt} =\mu^2+\lambda\nu,}\\[3mm]
{\displaystyle \frac{d\nu}{dt} =\nu^2+\lambda\mu.}
\end{cases}
\end{equation*}

Consider again the function $$y=f(x)=x(\log x-3)$$ for $e^2\leq
x<+\infty$, which is increasing and convex with range $-e^2\leq
y<+\infty$. Its inverse function $x=f^{-1}(y)$ is increasing and
concave on $-e^2\leq y<+\infty$. For each $t\geq0$, we consider
the set $K(t)$ of $3\times3$ symmetric matrices defined by the
inequalities: \begin{equation}
\lambda+\mu+\nu \geq-\frac{3}{1+t}, %\eqno(5.3.1)
\end{equation} and \begin{equation}
\nu(1+t)+f^{-1}((\lambda+\mu+\nu)(1+t))\geq0, %\eqno(5.3.2)
\end{equation} which is closed and convex (as we saw in the proof of Theorem
2.4.1). By the assumptions at $t=0$ and the advanced maximum
principle Theorem 2.3.5, we only need to check that the set $K(t)$
is preserved by the ODE system.

Since $R=\lambda+\mu+\nu$, we get from the ODE that
$$
\frac{dR}{dt}\geq\frac{2}{3}R^2\geq\frac{1}{3}R^2
$$
which implies that
$$
R\geq-\frac{3}{1+t},\ \ \ \text{for\ all }\ t\geq0.
$$
Thus the first inequality (5.3.1) is preserved. Note that the
second inequality (5.3.2) is automatically satisfied when $(-\nu)
\leq 3/(1+t)$. Now we compute from the ODE system,
\begin{displaymath}
\begin{split}
\frac{d}{dt}(\frac{R}{(-\nu)}-\log(-\nu))&  =\frac{1}{(-\nu)^2}
\left[(-\nu)\cdot\frac{dR}{dt}-(R+(-\nu))\frac{d(-\nu)}{dt}\right]\\[1mm]
&  =\frac{1}{(-\nu)^2}[(-\nu)^3+(-\nu)\mu^2+\lambda^2((-\nu)+\mu)
-\lambda\mu(\nu-\mu)]\\[1mm]
&  \geq(-\nu)\\[3mm]
&  \geq\frac{3}{(1+t)}\\[1mm]
&  \geq\frac{d}{dt}[\log(1+t)-3]
\end{split}
\end{displaymath}
whenever $R=(-\nu)[\log(-\nu)+\log(1+t)-3]$ and
$(-\nu)\geq3/(1+t)$. Thus the second inequality (5.3.2) is also
preserved under the system of ODE.

Therefore we have proved the theorem.
%$$\eqno \#$$ \vskip 0.3cm
\end{proof}

Denote by
$$
\hat{\rho}(t)=\max\{\operatorname{inj}\,(x,g_{ij}(t))\ |\ x\in M \}
$$
where $\operatorname{inj}\,(x,g_{ij}(t))$ is the injectivity radius of the
manifold $M$ at $x$ with respect to the metric $g_{ij}(t)$.

%\vskip0.2cm\noindent {\bf Definition 5.3.3}\ \

%CZ-SOURCE-LINE-11989
\begin{definition}[collapsed nonsingular normalized solution]
  \label{def:cz-collapsed-nonsingular-solution}
  \dcref{Cao--Zhu Definition 5.3.3}
  \group{Nonsingular Three-Dimensional Solutions}
  \source{cao-zhu}

We say a solution to the normalized Ricci flow is {\bf
collapsed}\index{collapsed} if there is a sequence of times
$t_k\rightarrow+\infty$ such that $\hat{\rho}(t_k)\rightarrow0$ as
$k\rightarrow+\infty$.
\end{definition}

When a nonsingular solution of the Ricci flow on $M$ is collapsed,
it follows from the work of Cheeger-Gromov \cite{ChG86, ChG90} or
Cheeger-Gromov-Fukaya \cite{CGF92} that the manifold $M$ has an
$\mathcal{F}$-structure and then its topology is completely
understood. Thus, in the following, we always assume our
nonsingular solutions are not collapsed.

Now suppose that we have a nonsingular solution which does not
collapse. Then for arbitrary sequence of times
$t_j\rightarrow\infty$, we can find a sequence of points $x_j$ and
some $\delta>0$ so that the injectivity radius of $M$ at $x_j$ in
the metric at time $t_j$ is at least $\delta$. Clearly the
Hamilton's compactness theorem (Theorem 4.1.5) also holds for the
normalized Ricci flow. Then by taking the $x_j$ as origins and the
$t_j$ as initial times, we can extract a convergent subsequence.
We call such a limit a {\bf noncollapsing
limit}\index{noncollapsing limit}. Of course the limit has also
finite volume. However the volume of the limit may be smaller than
the original one if the diameter goes to infinity.

The main result of this section is the following theorem of
Hamilton \cite{Ha99}.

%\vskip 0.2cm\noindent {\bf Theorem 5.3.4}\
%CZ-SOURCE-LINE-12020
\begin{theorem}[structure of noncollapsing nonsingular three-dimensional flows]
  \label{thm:cz-noncollapsing-nonsingular-flow-structure}
  \dcref{Cao--Zhu Theorem 5.3.4}
  \group{Nonsingular Three-Dimensional Solutions}
  \source{cao-zhu}

Let $g_{ij}(t)$, $0\leq t<+\infty$, be a noncollapsing nonsingular
solution of the normalized Ricci flow on a compact three-manifold
$M$. Then either
\begin{itemize}
\item[(i)] there exist a sequence of times $t_k\rightarrow+\infty$
and a sequence of diffeomorphisms $\varphi_k:\ M\rightarrow M$ so
that the pull-back of the metric $g_{ij}(t_k)$ by $\varphi_k$
converges in the $C^\infty$ topology to a metric on $M$ with
constant sectional curvature; or \item[(ii)] we can find a finite
collection of complete noncompact hyperbolic three-manifolds
$\mathcal{H}_1,\ldots,\mathcal{H}_m$ with finite volume, and for
all $t$ beyond some time $T<+\infty$ we can find compact subsets
$K_1,\ldots,K_m$ of $\mathcal{H}_1,\ldots,\mathcal{H}_m$
respectively obtained by truncating each cusp of the hyperbolic
manifolds along constant mean curvature torus of small area, and
diffeomorphisms $\varphi_l(t)$, $1\leq l\leq m$, of $K_l$ into $M$
so that as long as $t$ sufficiently large, the pull-back of the
solution metric $g_{ij}(t)$ by $\varphi_l(t)$ is as close as to
the hyperbolic metric as we like on the compact sets
$K_1,\ldots,K_m$; and moreover if we call the \textbf{exceptional
part} of $M$ those points where they are not in the image of any
$\varphi_l$, we can take the injectivity radii of the exceptional
part at everywhere as small as we like and the boundary tori of
each $K_l$ are \textbf{incompressible} in the sense that each
$\varphi_l$ injects $\pi_1(\partial K_l)$ into $\pi_1(M)$.
\index{exceptional part} \index{incompressible}
\end{itemize}
\end{theorem}

%\vskip 0.2cm\noindent {\bf Remark 5.3.5}\
%CZ-SOURCE-LINE-12051
\begin{remark}[geometry of the exceptional long-time region]
  \label{rem:cz-exceptional-long-time-region-geometry}
  \dcref{Cao--Zhu Remark 5.3.5}
  \group{Nonsingular Three-Dimensional Solutions}
  \source{cao-zhu}

The exceptional part has bounded curvature and arbitrarily small
injectivity radii everywhere as $t$ large enough. Moreover the
boundary of the exceptional part consists of a finite disjoint
union of tori with sufficiently small area and is convex. Then by
the work of Cheeger-Gromov \cite{ChG86}, \cite{ChG90} or
Cheeger-Gromov-Fukaya \cite{CGF92}, there exists an
$\mathcal{F}$-structure on the exceptional part. In particular,
the exceptional part is a graph manifold, which have been
topologically classified. Hence \textbf{any nonsingular solution
to the normalized Ricci flow is geometrizable} in the sense of
Thurston (see the last section of Chapter 7 for details).
\end{remark}

The rest of this section is devoted to the proof of Theorem 5.3.4.
We now present the proof given by Hamilton \cite{Ha99} and will
divide his arguments in \cite{Ha99} into the following three
parts.

\vskip 0.5cm {\bf Part I: Subsequence Convergence} \vskip 0.2cm

According to Lemma 5.1.1, the scalar curvature of the normalized
flow evolves by the equation
\begin{align}
\frac{\partial }{\partial t}R
&  =\Delta R+2|Ric|^2-\frac{2}{3}rR\\
&  = \Delta R+2|\stackrel{\circ}{Ric}|^2+\frac{2}{3}R(R-r)\nonumber
\end{align}  %\eqno(5.3.3)
where $\stackrel{\circ}{Ric}$ is the traceless part of the Ricci
tensor. As before, we denote by $R_{\min}(t)=\min_{x\in M}R(x,t)$.
It then follows from the maximum principle that \begin{equation}
\frac{d}{dt}R_{\min}\geq\frac{2}{3}R_{\min}(R_{\min}-r),
%\eqno(5.3.4)
\end{equation} which implies that if $R_{\min}\leq0$ it must be
nondecreasing, and if $R_{\min}\geq0$ it cannot go negative again.
We can then divide the noncollapsing solutions of the normalized
Ricci flow into three cases.

\smallskip
{\it Case} (1): $R_{\min}(t)>0$ for some $t>0$;

\smallskip
{\it Case} (2): $R_{\min}(t)\leq0$ for all $t\in[0,+\infty)$ and
$\lim\limits_{t\rightarrow+\infty}R_{\min}(t)=0$;

\smallskip
{\it Case} (3): $R_{\min}(t)\leq0$ for all $t\in[0,+\infty)$ and
$\lim\limits_{t\rightarrow+\infty}R_{\min}(t)<0$.

\smallskip
Let us first consider Case (1). In this case the maximal time
interval $[0,T)$ of the corresponding solution of the unnormalized
flow is finite, since the unnormalized scalar curvature
$\tilde{R}$ satisfies
\begin{displaymath}
\begin{split}
\frac{\partial}{\partial t}\tilde{R}
&  =\Delta\tilde{R}+2|\tilde{R}ic|^2\\
&  \geq\Delta\tilde{R}+\frac{2}{3}\tilde{R}^2
\end{split}
\end{displaymath}
which implies that the curvature of the unnormalized solution
blows up in finite time. Without loss of generality, we may assume
that for the initial metric at $t=0$, the eigenvalues
$\tilde{\lambda}\geq\tilde{\mu}\geq\tilde{\nu}$ of the curvature
operator are bounded below by $\tilde{\nu}\geq-1$. It follows from
Theorem 5.3.2 that the pinching estimate
$$
\tilde{R}\geq(-\tilde{\nu})[\log(-\tilde{\nu})+\log(1+t)-3]
$$
holds whenever $\tilde{\nu}<0$. This shows that when the
unnormalized curvature big, the negative ones are not nearly as
large as the positive ones. Note that the unnormalized curvature
becomes unbounded in finite time. Thus when we rescale the
unnormalized flow to the normalized flow, the scaling factor must
go to infinity. In the nonsingular case the rescaled positive
curvature stay finite, so the rescaled negative curvature (if any)
go to zero. Hence we can take a noncollapsing limit for the
nonsingular solution of the normalized flow so that it has
nonnegative sectional curvature.

Since the volume of the limit is finite, it follows from a result
of Calabi and Yau (cf. \cite{ScY}) that the limit must be compact
and the limiting manifold is the original one. Then by the strong
maximum principle as in the proof of Theorem 5.2.3 (i), either the
limit is flat, or it is a compact metric quotient of the product
of a positively curved surface $\Sigma^2$ with $\mathbb{R}$, or it
has strictly positive curvature. By the work of Schoen-Yau
\cite{ScY79}, a flat three-manifold cannot have a metric of
positive scalar curvature, but our manifold does in Case (1). This
rules out the possibility of a flat limit. Clearly the limit is
also a nonsingular solution to the normalized Ricci flow. Note
that the curvature of the surface $\Sigma^2$ has a positive lower
bound and is compact since it comes from the lifting of the
compact limiting manifold. From Theorem 5.1.11, we see the metric
of the two-dimensional factor $\Sigma^2$ converges to the round
two-sphere $\mathbb{S}^2$ in the normalized Ricci flow. Note also
that the normalized factors in two-dimension and three-dimension
are different. This implies that the compact quotient of the
product $\Sigma^2 \times \mathbb{R}$ cannot be nonsingular, which
is also ruled out for the limit. Thus the limit must have strictly
positive sectional curvature. Since the convergence takes place
everywhere for the compact limit, it follows that as $t$ large
enough the original nonsingular solution has strictly positive
sectional curvature. This in turn shows that the corresponding
unnormalized flow has strictly positive sectional curvature after
some finite time. Then in views of the proof of Theorem 5.2.1, in
particular the pinching estimate in Proposition 5.2.5, the limit
has constant Ricci curvature and then constant sectional curvature
for three-manifolds. This finishes the proof in Case (1).

We next consider Case (2). In this case we only need to show that
we can take a noncollapsing limit which has nonnegative sectional
curvature. Indeed, if this is true, then as in the previous case,
the limit is compact and either it is flat, or it splits as a
product (or a quotient of a product) of a positively curved $S^2$
with a circle $S^1$, or it has strictly positive curvature. But
the assumption $R_{\min}(t)\leq0$ for all times $t\geq0$ in this
case implies the limit must be flat.

Let us consider the corresponding unnormalized flow
$\tilde{g}_{ij}(t)$ associated to the noncollapsing nonsingular
solution. The pinching estimate in Theorem 5.3.2 tells us that we
may assume the unnormalized flow $\tilde{g}_{ij}(t)$ exists for
all times $0\leq t<+\infty$, for otherwise, the scaling factor
approaches infinity as in the previous case which implies the
limit has nonnegative sectional curvature. The volume
$\tilde{V}(t)$ of the unnormalized solution $\tilde{g}_{ij}(t)$
now changes. We divide the discussion into three subcases.

Subcase (2.1): there is a sequence of times
$\tilde{t}_k\rightarrow+\infty$ such that
$\tilde{V}(\tilde{t}_k)\rightarrow+\infty$;

Subcase (2.2): there is a sequence of times
$\tilde{t}_k\rightarrow+\infty$ such that
$\tilde{V}(\tilde{t}_k)\rightarrow0$;

Subcase (2.3): there exist two positive constants $C_1$, $C_2$
such that $C_1\leq\tilde{V}(t)\leq C_2$ for all $0\leq t<+\infty$.

For Subcase (2.1), because
$$
\frac{d\tilde{V}}{dt}=-r\tilde{V}
$$
we have
$$
\log\frac{\tilde{V}(\tilde{t}_k)}{\tilde{V}(0)}
=-\int_0^{\tilde{t}_k}r(t)dt\rightarrow+\infty,\; \text{ as }\;
k\rightarrow+\infty,
$$
which implies that there exists another sequence of times, still
denoted by $\tilde{t}_k$, such that
$\tilde{t}_k\rightarrow+\infty$ and $r(\tilde{t}_k)\leq0$. Let
$t_k$ be the corresponding times for the normalized flow. Thus
there holds for the normalized flow
$$
r(t_k)\rightarrow0,\quad\text{as }\; k\rightarrow\infty,
$$
since $0\geq r(t_k)\geq R_{\min}(t_k)\rightarrow0$ as
$k\rightarrow+\infty$. Then
$$
\int_M(R-R_{\min})d\mu(t_k)=(r(t_k)-R_{\min}(t_k))V\rightarrow0,
\quad \text{as }\; k\rightarrow\infty.
$$
As we take a noncollapsing limit along the time sequence $t_k$, we
get
$$
\int_{M^\infty}Rd\mu^\infty=0
$$
for the limit of the normalized solutions at the new time $t=0$.
But $R\geq0$ for the limit because
$\lim\limits_{t\rightarrow+\infty}R_{\min}(t)=0$ for the
nonsingular solution. So $R=0$ at $t=0$ for the limit. Since the
limit flow exists for $-\infty<t<+\infty$ and the scalar curvature
of the limit flow evolves by
$$
\frac{\partial}{\partial t}R=\Delta
R+2|\Ric|^2-\frac{2}{3}r^\infty R,\quad t\in(-\infty,+\infty)
$$
where $r^\infty$ is the limit of the function $r(t)$ by
translating the times $t_k$ as the new time $t=0$. It follows from
the strong maximum principle that
$$
R\equiv0, \quad \text{on }\; M^\infty\times(-\infty,+\infty).
$$
This in turn implies, in view of the above evolution equation,
that
$$
\Ric\equiv0, \quad \text{on }\; M^\infty\times(-\infty,+\infty).
$$
Hence this limit must be flat. Since the limit $M^\infty$ is
complete and has finite volume, the flat manifold $M^\infty$ must
be compact. Thus the underlying manifold $M^\infty$ must agree
with the original $M$ (as a topological manifold). This says that
the limit was taken on $M$.

For Subcase (2.2), we may assume as before that for the initial
metric at $t=0$ of the unnormalized flow $\tilde{g}_{ij}(t)$, the
eigenvalues $\tilde{\lambda}\geq\tilde{\mu}\geq\tilde{\nu}$ of the
curvature operator satisfy $\tilde{\nu}\geq-1$. It then follows
from Theorem 5.3.2 that
$$
\tilde{R}\geq(-\tilde{\nu})[\log(-\tilde{\nu})+\log(1+t)-3], \quad
\text{for all }\; t\geq0
$$
whenever $\tilde{\nu}<0$.

Let $t_k$ be the sequence of times in the normalized flow which
corresponds to the sequence of times $\tilde{t}_k$. Take a
noncollapsing limit for the normalized flow along the times $t_k$.
Since $\tilde{V}(\tilde{t}_k)\rightarrow 0$, the normalized
curvatures at the times $t_k$ are reduced by multiplying the
factor $(\tilde{V}(\tilde{t}_k))^\frac{2}{3}$. We claim the
noncollapsing limit has nonnegative sectional curvature. Indeed if
the maximum value of $(-\tilde{\nu})$ at the time $\tilde{t}_k$
does not go to infinity, the normalized eigenvalue $-\nu$ at the
corresponding time $t_k$ must get rescaled to tend to zero; while
if the maximum value of $(-\tilde{\nu})$ at the time $\tilde{t}_k$
does go to infinity, the maximum value of $\tilde{R}$ at
$\tilde{t}_k$ will go to infinity even faster from the pinching
estimate, and when we normalize to keep the normalized scalar
curvature $R$ bounded at the time $t_k$ so the normalized
$(-{\nu})$ at the time $t_k$ will go to zero. Thus in either case
the noncollapsing limit has nonnegative sectional curvature at the
initial time $t=0$ and then has nonnegative sectional curvature
for all times $t\ge 0$.

For Subcase (2.3), normalizing the flow only changes quantities in
a bounded way. As before we have the pinching estimate
$$R\ge (-\nu)[\log(-\nu)+\log(1+t)-C]$$
for the normalized Ricci flow, where $C$ is a positive constant
depending only on the constants $C_1$, $C_2$ in the assumption of
Subcase (2.3). If $$(-\nu)\le \frac{A}{1+t}$$ for any fixed
positive constant $A$, then $(-\nu)\rightarrow 0$ as $t\rightarrow
+\infty$ and we can take a noncollapsing limit which has
nonnegative sectional curvature. On the other hand if we can pick
a sequence of times $t_k\rightarrow\infty$ and points $x_k$ where
$(-\nu)(x_k,t_k)=\max\limits_{x\in M}(-\nu)(x,t_k)$ satisfies
$$
(-\nu)(x_k,t_k)(1+t_k)\rightarrow +\infty,\quad \text{as } \;
k\rightarrow +\infty,
$$
then from the pinching estimate, we have
$$
\frac{R(x_k,t_k)}{(-\nu)(x_k,t_k)}\rightarrow +\infty,\quad
\text{as}\quad k\rightarrow +\infty.
$$
But $R(x_k,t_k)$ are uniformly bounded since normalizing the flow
only changes quantities in bounded way. This shows $\sup
(-\nu)(\cdot,t_k)\rightarrow 0$ as $k\rightarrow+\infty$. Thus we
can take a noncollapsing limit along $t_k$ which has nonnegative
sectional curvature. Hence we have completed the proof of Case
(2).

We now come to the most interesting Case (3) where $R_{min}$
increases monotonically to a limit strictly less than zero. By
scaling we can assume $R_{min}(t)\rightarrow -6$ as $t\rightarrow
+\infty$.

%\vskip 0.2cm \noindent{\bf Lemma 5.3.6} \ \emph{
%CZ-SOURCE-LINE-12312
\begin{lemma} [hyperbolicity of noncollapsing limits]
  \label{lem:cz-noncollapsing-limits-hyperbolic}
  \dcref{Cao--Zhu Lemma 5.3.6}
  \group{Nonsingular Three-Dimensional Solutions}
  \source{cao-zhu}

In Case $(3)$ where $R_{\min}\rightarrow-6$ as $t\rightarrow
+\infty$, all noncollapsing limit are hyperbolic with constant
sectional curvature $-1.$
\end{lemma}

%\noindent {\bf Proof.}\
\begin{proof}
By (5.3.4) and the fact $R_{\min}(t)\le-6$, we have
$$
\frac{d}{dt}R_{\min}(t)\ge 4(r(t)-R_{\min}(t))
$$
and
$$
\int_0^\infty(r(t)-R_{\min}(t))dt<+\infty.
$$
Since $r(t)-R_{\min}(t)\ge 0$ and $R_{min}(t)\rightarrow-6$ as
$t\rightarrow+\infty$, it follows that the function $r(t)$ has the
limit
$$
r=-6,
$$
for any convergent subsequence. And since
$$
\int_M(R-R_{\min}(t))d\mu=(r(t)-R_{\min}(t))\cdot V,
$$
it then follows that
$$
R\equiv-6\quad \text{for the limit}.
$$
The limit still has the following evolution equation for the
limiting scalar curvature
$$
\frac{\partial}{\partial t}R=\Delta R+2| {\mathop {\rm
Ric}\limits^\circ}|^2+\frac{2}{3}R(R-r).
$$
Since $R\equiv r\equiv -6$ in space and time for the limit, it
follows directly that $|{\mathop {\rm Ric}\limits^ \circ}|\equiv
0$ for the limit. Thus the limit metric has $\lambda=\mu=\nu=-2$,
so it has constant sectional curvature $-1$ as desired.
%$$\eqno \#$$ \vskip 0.2cm
\end{proof}

If in the discussion above there exists a compact noncollapsing
limit, then we know that the underlying manifold $M$ is compact
and we fall into the conclusion of Theorem 5.3.4(i) for the
constant negative sectional curvature limit. Thus it remains to
show when every noncollapsing limit is a complete noncompact
hyperbolic manifold with finite volume, we have conclusion (ii) in
Theorem 5.3.4.

Now we first want to find a finite collection of persistent
complete noncompact hyperbolic manifolds as stated in Theorem
5.3.4 (ii).
%\pagebreak

\bigskip
{\bf Part II: Persistence of Hyperbolic Pieces}

\smallskip
We begin with the definition of the topology of $C^\infty$
convergence on compact sets for maps $F:M\rightarrow N$ of one
Riemannian manifold to another. For any compact set
$K\subset\subset M$ and any two maps $F,G:M\rightarrow N$, we
define
$$
d_K(F,G)=\sup_{x\in K}d(F(x),G(x))
$$
where $d(y,z)$ is the geodesic distance from $y$ to $z$ on $N$.
This gives the $C^0_{loc}$ topology for maps between $M$ and $N$.
To define $C^k_{loc}$ topology for any positive integer $k\ge1$,
we consider the \mbox{\boldmath{$k$}}-{\bf jet
space}\index{$k$-jet!space} $J^kM$ of a manifold $M$ which is the
collection of all
$$
(x,J^1,J^2,\ldots,J^k)
$$
where $x$ is a point on $M$ and $J^i$ is a tangent vector for
$1\le i\le k$ defined by the $i^{th}$ covariant derivative
$J^i=\nabla^i_\frac{\partial}{\partial t}\gamma(0)$ for a path
$\gamma$ passing through the point $x$ with $\gamma(0)=x$. A
smooth map $F:M\rightarrow N$ induces a map $$J^kF:\quad
J^kM\rightarrow J^kN$$ defined by
$$
J^kF(x,J^1,\ldots,J^k)=(F(x),\nabla_\frac{\partial}{\partial t}
(F(\gamma))(0),\ldots, \nabla^k_\frac{\partial}{\partial
t}(F(\gamma))(0))
$$
where $\gamma$ is a path passing through the point $x$ with
$J^i=\nabla^i_\frac{\partial}{\partial t}\gamma(0),1\le i\le k.$

Define the {\bf ${\mathbf k}$-jet
distance}\index{$k$-jet!distance} between $F$ and $G$ on a compact
set $K\subset\subset M$ by
$$
d_{C^k(K)}(F,G)=d_{BJ^kK}(J^kF,J^kG)
$$
where $BJ^kK$ consists of all $k$-jets $(x,J^1,\ldots,J^k)$ with
$x\in K$ and
$$
|J^1|^2+|J^2|^2+\cdots+|J^k|^2\le 1.
$$
Then the convergence in the metric $d_{C^k(K)}$ for all positive
integers $k$  and all compact sets $K$ defines the topology of
$C^\infty$ convergence on compact sets for the space of maps.

We will need the following \textbf{Mostow type rigidity}
\index{Mostow type rigidity} result (cf. Corollary 8.3 of
\cite{Ha99}).

%\vskip 0.2cm \noindent{\bf Lemma 5.3.7} \ \emph{
%CZ-SOURCE-LINE-12423
\begin{lemma}[stability of finite-volume hyperbolic embeddings]
  \label{lem:cz-finite-volume-hyperbolic-embedding-stability}
  \dcref{Cao--Zhu Lemma 5.3.7}
  \group{Nonsingular Three-Dimensional Solutions}
  \source{cao-zhu}
 For any complete noncompact
hyperbolic three-manifold $\mathcal{H}$ with finite volume with
metric $h$, we can find a compact set $\mathcal{K}$ of
$\mathcal{H}$ such that for every integer $k$ and every
$\varepsilon>0$, there exist an integer $q$ and a $\delta>0$ with
the following property: if $F$ is a diffeomorphism of
$\mathcal{K}$ into another complete noncompact hyperbolic
three-manifold $\tilde{\mathcal{H}}$ with no fewer cusps $($than
$\mathcal{H}),$ finite volume with metric $\tilde{h}$ such that
$$
\|F^*\tilde{h}-h\|_{C^q(\mathcal{K})}<\delta
$$
then there exists an isometry $I$ of $\mathcal{H}$ to
$\tilde{\mathcal{H}}$ such that
$$
d_{C^k(\mathcal{K})}(F,I)<\varepsilon.
$$
\end{lemma}

%\vskip 0.1cm\noindent{\bf{Proof.}}
\begin{proof} This version of Mostow rigidity is given by Hamilton
(cf. section 8 of \cite{Ha99}). The following argument is in part
based on the Editors' notes in p.323-324 of \cite{CCCY}.

First we claim that $\mathcal{H}$ is isometric to
$\mathcal{\tilde{H}}$ for an appropriate choice of compact set
$\mathcal{K}$, positive integer $q$ and positive number $\delta$.
Let $l: \mathcal{H}\rightarrow \mathbb{R}$ be a function defined
at each point by the length of the shortest non-contractible loop
starting and ending at this point. Denote the Margulis constant by
$\mu$. Then by Margulis lemma (see for example \cite{Gro79} or
\cite{KM}), for any $0<\varepsilon_0<\frac{1}{2}\mu$, the set
$l^{-1}([0,\varepsilon_0])\subset \mathcal{H}$ consists of
finitely many components and each of these components is isometric
to a cusp or to a tube. Topologically, a tube is just a solid
torus. Let $\varepsilon_0$ be even smaller than one half of the
minimum of the lengths of the all closed geodesics on the tubes.
Then $l^{-1}([0,\varepsilon_0])$ consists of finite number of
cusps. Set $\mathcal{K}_0 = l^{-1}([\varepsilon_0,\infty))$. The
boundary of $\mathcal{K}_0$ consists of flat tori with constant
mean curvatures. Note that each embedded torus in a complete
hyperbolic three-manifold with finite volume either bounds a solid
torus or is isotopic to a standard torus in a cusp. The
diffeomorphism $F$ implies the boundary $F(\partial\mathcal{K}_0)$
are embedded tori. If one of components bounds a solid torus, then
as $\delta$ sufficiently small and $q$ sufficiently large,
$\mathcal{\tilde{H}}$ would have fewer cusps than $\mathcal{H}$,
which contradicts with our assumption. Consequently,
$\mathcal{\tilde{H}}$ is diffeomorphic to
$F({\overset{o}{\mathcal{K}_0}})$. Here
${\overset{o}{\mathcal{K}_0}}$ is the interior of the set
$\mathcal{K}_0$.  Since $\mathcal{H}$ is diffeomorphic to
${\overset{o}{\mathcal{K}_0}}$, $\mathcal{H}$ is diffeomorphic to
$\mathcal{\tilde{H}}$.  Hence by Mostow's rigidity theorem (see
\cite{Mos} and \cite{Pr}), $\mathcal{H}$ is isometric to
$\mathcal{\tilde{H}}$.

So we can assume $\mathcal{\tilde{H}}=\mathcal{H}$.  For
$\mathcal{K} = \mathcal{K}_0$, we argue by contradiction. Suppose
there is some $k>0$ and $\varepsilon>0$ so that there exist
sequences of integers $q_j\rightarrow\infty$, $\delta_j\rightarrow
0^+$ and diffeomorphisms $F_j$ mapping $\mathcal{K}$ into
$\mathcal{H}$ with
$$
  \|F_j^{*}h-h\|_{C^{q_j}(\mathcal{K})}<\delta_j
$$
 and
$$
 d_{C^{k}(\mathcal{K})}(F_j,I)\geq \varepsilon
$$
for all isometries $I$ of $\mathcal{H}$ to itself. We can extract
a subsequence of $F_j$ convergent to a map $F_{\infty}$ with
$F_{\infty}^{*}h=h$ on $\mathcal{K}$.

We need to check that $F_{\infty}$ is still a diffeomorphism on
$\mathcal{K}$. Since $F_{\infty}$ is a local diffeomorphism and is
the limit of diffeomorphisms, we can find an inverse of
$F_{\infty}$ on $F_{\infty}(\overset{o}{\mathcal{K}})$. So
$F_{\infty}$ is a diffeomorphism on $\overset{o}{\mathcal{K}}$. We
claim the image of the boundary can not touch the image of the
interior. Indeed, if $F_{\infty}(x_{1})=F_{\infty}(x_{2})$ with
$x_{1}\in \partial \mathcal{K} $ and $x_{2}\in
\overset{o}{\mathcal{K}} $, then we can find $x_{3}\in
\overset{o}{\mathcal{K}}$ near $x_{1}$ and $x_{4}\in
\overset{o}{\mathcal{K}}$ near $x_{2}$ with
$F_{\infty}(x_{3})=F_{\infty}(x_{4})$, since $F_{\infty}$ is a
local diffeomorphism. This contradicts with the fact that
$F_{\infty}$ is a diffeomorphism on $\overset{o}{\mathcal{K}}$.
This proves our claim. Hence, the only possible overlap is at the
boundary. But the image $F_{\infty}(\partial \mathcal{K})$ is
strictly concave, this prevents the boundary from touching itself.
We conclude that the mapping $F_{\infty}$ is a diffeomorphism on
$\mathcal{K}$, hence an isometry.

To extend $F_{\infty}$ to a global isometry, we argue as follows.
For each truncated cusp end of $\mathcal{K}$,  the area of
constant mean curvature flat torus is strictly decreasing. Since
$F_{\infty}$ takes each such torus to another of the same area, we
see that $F_{\infty}$ takes the foliation of an end by constant
mean curvature flat tori to another such foliation. So
$F_{\infty}$ takes cusps to cusps and preserves their foliations.
Note that the isometric type of a cusp is just the isometric type
of the torus, more precisely, let
$(N,dr^{2}+e^{-2r}g_{\mathcal{V}})$ be a cusp (where
$g_{\mathcal{V}}$ is the flat metric on the torus $\mathcal{V}$),
$0<a<b$ are two constants, any isometry of $N\cap l^{-1}[a,b]$ to
itself is just an isometry of $\mathcal{V}$. Hence the isometry
$F_{\infty}$ can be extended to the whole cusps. This gives a
global isometry $I$ contradicting our assumption when $j$ large
enough.

The proof of the Lemma 5.3.7 is completed.
%$$\eqno \#$$
\end{proof}

In order to obtain the persistent hyperbolic pieces stated in
Theorem 5.3.4 (ii), we will need to use a special parametrization
given by harmonic maps.

%\vskip 0.2cm \noindent {\bf Lemma 5.3.8}  \ \emph{
%CZ-SOURCE-LINE-12543
\begin{lemma} [metric perturbations with concave boundary]
  \label{lem:cz-concave-boundary-metric-perturbations}
  \dcref{Cao--Zhu Lemma 5.3.8}
  \group{Nonsingular Three-Dimensional Solutions}
  \source{cao-zhu}

Let\; $(X,\,g)$\; be\; a\; compact\; Riemannian\; manifold\; with
strictly negative Ricci curvature and with strictly concave
boundary. Then there are positive integer $l_0$ and small number
$\varepsilon_0>0$ such that for each positive integer $l\ge l_0$ and
positive number $\varepsilon\le\varepsilon_0$ we can find positive
integer $q$ and positive number $\delta>0$ such that for every
metric $\tilde{g}$ on $X$ with $||\tilde{g}-g||_{C^q(X)}\le \delta$
we can find a unique diffeomorphism $F$ of $X$ to itself so that
\begin{itemize}
\item[(a)] $F: (X,g)\rightarrow (X,\tilde{g})$ is harmonic,
\item[(b)] $F$ takes the boundary $\partial X$ to itself and
satisfies the {\bf free boundary condition}\index{free boundary
condition} that the normal derivative $\nabla_NF$ of $F$ at the
boundary is normal to the boundary, \item[(c)]
$d_{C^l(X)}(F,Id)<\varepsilon$, where $Id$ is the identity map.
\end{itemize}
\end{lemma}

%\vskip 0.1cm\noindent {\bf Proof.} \
\begin{proof}
  \uses{lem:cz-hyperbolic-cusp-length-bounded}
The following argument is adapted from Hamilton's paper
\cite{Ha99} and the Editors' note on p.325 of \cite{CCCY}. Let
$\Phi(X,\partial X)$ be the space of maps of $X$ to itself which
take $\partial X$ to itself. Then $\Phi(X,\partial X)$ is a Banach
manifold and the tangent space to $\Phi(X,\partial X)$ at the
identity is the space of vector fields
$V=V^i\frac{\partial}{\partial x^i}$ tangent to the boundary.
Consider the map sending $F\in \Phi(X,\partial X)$ to the pair
$\{\Delta F,(\nabla_NF)_{//}\}$ consisting of the harmonic map
Laplacian and the tangential component (in the target) of the
normal derivative of $F$ at the boundary. By using the inverse
function theorem, we only need to check that the derivative of
this map is an isomorphism at the identity with $\tilde{g}=g$.

Let $\{x^i\}_{i=1,\ldots,n}$ be a local coordinates of $(X,g)$ and
$\{y^\alpha\}_{\alpha=1,\ldots,n}$ be a local coordinates of
$(X,\tilde{g})$. The harmonic map Laplacian of $F:
(X,g)\rightarrow (X,\tilde{g})$ is given in local coordinates by
$$
(\Delta F)^\alpha
=\Delta(F^\alpha)+g^{ij}(\tilde{\Gamma}^\alpha_{\beta\gamma}\circ
F) \frac{\partial F^\beta}{\partial x^i}\frac{\partial
F^\gamma}{\partial x^j}
$$
where $\Delta(F^\alpha)$ is the Laplacian of the function
$F^\alpha$ on $X$ and $\tilde{\Gamma}^\alpha_{\beta\gamma}$ is the
connection of $\tilde{g}$. Let $F$ be a one-parameter family with
$F|_{s=0}=Id$ and $\frac{dF}{ds}|_{s=0}=V$, a smooth vector field
on $X$ tangent to the boundary (with respect to $g$). At an
arbitrary given point $x\in X$, we choose the coordinates
$\{x^i\}_{i=1,\ldots,n}$ so that $\Gamma^i_{jk}(x)=0.$ We compute
at the point $x$ with $\tilde{g}=g$,
$$
\frac{d}{ds}{\Big|_{s=0}}(\Delta F)^\alpha
=\Delta(V^\alpha)+g^{ij}\left(\frac{\partial}{\partial
x^k}\Gamma^\alpha_{ij}\right)V^k.
$$
Since
$$(\nabla_iV)^{\alpha} =
\nabla_iV^{\alpha} + (\Gamma^{\alpha}_{i\beta}\circ F)V^{\beta},
$$
we have, at $s=0$ and the point $x$,
$$
(\Delta V)^\alpha=\Delta(V^\alpha)+g^{ij}\frac{\partial}{\partial
x^i}\Gamma^\alpha_{jk}V^k.
$$
Thus we obtain
\begin{align}
\frac{d}{ds}{|_{s=0}}(\Delta F)^\alpha&  = (\Delta
V)^\alpha+g^{ij}\left(\frac{\partial}{\partial
x^k}\Gamma^\alpha_{ij}-\frac{\partial}{\partial
x^i}\Gamma^\alpha_{jk}\right)V^k \\
&  = (\Delta V)^\alpha+g^{\alpha i}R_{ik}V^k.\nonumber
\end{align} %\eqno(5.3.5)
Since
$$
(\nabla_NF)(F^{-1}(x)) =N^i(F^{-1}(x))\frac{\partial F^j}{\partial
x^i}(F^{-1}(x)) \frac{\partial}{\partial x^j}(x)\quad\text{on
}\;\partial X,
$$
we have
\begin{align}
\frac{d}{ds}|_{s=0}\{(\nabla_NF)_{//}\}
&  = \frac{d}{ds}|_{s=0}(\nabla_NF -\langle \nabla_NF,N \rangle N)\\
&  = \frac{d}{ds}|_{s=0}(\nabla_NF)-\left\langle
\frac{d}{ds}|_{s=0}\nabla_NF,N \right\rangle N  \nonumber\\
&\quad - \langle \nabla_NF,\nabla_VN \rangle N |_{s=0} - \langle
\nabla_NF,N
\rangle \nabla_VN |_{s=0} \nonumber\\
&  = \left(\frac{d}{ds}|_{s=0}\nabla_NF\right)_{//} -\nabla_VN \nonumber\\
&  = \left(-V(N^i)\frac{\partial}{\partial
x^i}+N(V^j)\frac{\partial}{\partial x^j}\right)_{//} - II(V) \nonumber\\
&  = [N,V]_{//} - II(V) \nonumber\\
&  = (\nabla_NV)_{//} - 2II(V) \nonumber
\end{align} %\eqno(5.3.6)
where $II$ is the second fundamental form of the boundary (as an
automorphism of $T(\partial X)$). Thus by (5.3.5) and (5.3.6), the
kernel of the map sending $F\in \Phi(X,\partial X)$ to the pair
$\{\Delta F,(\nabla_NF)_{//}\}$ is the space of solutions of
elliptic boundary value problem \begin{equation}
\begin{cases}
\Delta V+\Ric(V)=0 &\quad \text{on }\;X\\
V_{\bot}=0,&\quad \text{at }\;\partial X,\\
(\nabla_NV)_{//}-2II(V)=0,&\quad \text{at } \;\partial X,
\end{cases} %\eqno(5.3.7)
\end{equation} where $V_\bot$ is the normal component of $V$.

Now using these equations and integrating by parts gives
$$
\int\int\limits_X|\nabla V|^2
=\int\int\limits_X\Ric(V,V)+2\int\limits_{\partial X}II(V,V).
$$
Since $Rc<0$ and $II<0$ we conclude that the kernel is trivial.
Clearly this elliptic boundary value is self-adjoint because of
the free boundary condition. Thus the cokernel is trivial also.
This proves the lemma.
%$$\eqno \#$$ \vskip 0.3cm
\end{proof}

Now we can prove the persistence of hyperbolic pieces. Let
$g_{ij}(t),\;0\le t< +\infty$, be a noncollapsing nonsingular
solution of the normalized Ricci flow on a compact three-manifold
$M$. Assume that any noncollapsing limit of the nonsingular
solution is a complete noncompact hyperbolic three-manifold with
finite volume. Consider all the possible hyperbolic limits of the
given nonsingular solution, and among them choose one such
complete noncompact hyperbolic three-manifold $\mathcal{H}$ with
the least possible number of cusps. In particular, we can find a
sequence of times $t_k\rightarrow +\infty$ and a sequence of
points $P_k$ on $M$ such that the marked three-manifolds
$(M,g_{ij}(t_k),P_k)$ converge in the $C^\infty_{loc}$ topology to
$\mathcal{H}$ with hyperbolic metric $h_{ij}$ and marked point
$P\in \mathcal{H}$.

For any small enough $a>0$ we can truncate each cusp of
$\mathcal{H}$ along a constant mean curvature torus of area $a$
which is uniquely determined; the remainder we denote by
$\mathcal{H}_a$. Clearly as $a\rightarrow 0$ the $\mathcal{H}_a$
exhaust $\mathcal{H}$. Pick a sufficiently small number $a>0$ to
truncate cusps so that Lemma 5.3.7 is applicable for the compact
set $\mathcal{K} = \mathcal{H}_a$. Choose an integer $l_0$ large
enough and an $\varepsilon_0$ sufficiently small to guarantee from
Lemma 5.3.8 the uniqueness of the identity map $Id$ among maps
close to $Id$ as a harmonic map $F$ from $\mathcal{H}_a$ to itself
with taking $\partial\mathcal{H}_a$ to itself, with the normal
derivative of $F$ at the boundary of the domain normal to the
boundary of the target, and with $d_{C^{l_0}(\mathcal{H}_a)}(F,Id)
< \varepsilon_0$. Then choose positive integer $q_0$ and small
number $\delta_0>0$ from Lemma 5.3.7 such that if $\tilde{F}$ is a
diffeomorphism of $\mathcal{H}_a$ into another complete noncompact
hyperbolic three-manifold $\tilde{\mathcal{H}}$ with no fewer
cusps (than $\mathcal{H}$), finite volume with metric
$\tilde{h}_{ij}$ satisfying
$$
||\tilde{F}^*\tilde{h}_{ij}-h_{ij}||_{C^{q_0}(\mathcal{H}_a)} \le
\delta_0,
$$
then there exists an isometry $I$ of $\mathcal{H}$ to
$\tilde{\mathcal{H}}$ such that \begin{equation}
d_{C^{l_0}(\mathcal{H}_a)}(\tilde{F},I)<\varepsilon_0. %\eqno(5.3.8)
\end{equation} And we further require $q_0$ and $\delta_0$ from Lemma 5.3.8
to guarantee the existence of harmonic diffeomorphism from
$(\mathcal{H}_a,\tilde{g}_{ij})$ to $(\mathcal{H}_a,h_{ij})$ for
any metric $\tilde{g}_{ij}$ on $\mathcal{H}_a$ with
$||\tilde{g}_{ij}-h_{ij}||_{C^{q_0}(\mathcal{H}_a)}\le\delta_0.$

By definition, there exist a sequence of exhausting compact sets
$U_k$ of $\mathcal{H}$ (each $U_k \supset \mathcal{H}_a$) and a
sequence of diffeomorphisms $F_k$ from $U_k$ into $M$ such that
$F_k(P)=P_k$ and
$||F^*_kg_{ij}(t_k)-h_{ij}||_{C^m(U_k)}\rightarrow 0$ as
$k\rightarrow+\infty$ for all positive integers $m$. Note that
$\partial \mathcal{H}_a$ is strictly concave and we can foliate a
neighborhood of $\partial \mathcal{H}_a$ with constant mean
curvature hypersurfaces where the area $a$ has a nonzero gradient.
As the approximating maps
$F_k:\;(U_k,h_{ij})\rightarrow(M,g_{ij}(t_k))$ are close enough to
isometries on this collar of $\partial \mathcal{H}_a$, the metrics
$g_{ij}(t_k)$ on $M$ will also admit a unique constant mean
curvature hypersurface with the same area $a$ near $F_k(\partial
\mathcal{H}_a)(\subset M)$ by the inverse function theorem. Thus
we can change the map $F_k$ by an amount which goes to zero as
$k\rightarrow\infty$ so that now $F_k(\partial \mathcal{H}_a)$ has
constant mean curvature with the area $a$. Furthermore, by
applying Lemma 5.3.8 we can again change $F_k$ by an amount which
goes to zero as $k\rightarrow\infty$ so as to make $F_k$ a
harmonic diffeomorphism and take $\partial \mathcal{H}_a$ to the
constant mean curvature hypersurface $F_k(\partial \mathcal{H}_a)$
and also satisfy the free boundary condition that the normal
derivative of $F_k$ at the boundary of the domain is normal to the
boundary of the target. Hence for arbitrarily given positive
integer $q\ge q_0$ and positive number $\delta<\delta_0$, there
exists a positive integer $k_0$ such that for the modified
harmonic diffeomorphism $F_k$, when $k\ge k_0$,
$$
||F^*_kg_{ij}(t_k)-h_{ij}||_{C^q(\mathcal{H}_a)}<\delta.
$$

For each fixed $k\ge k_0$, by the implicit function theorem we can
first find a constant mean curvature hypersurface near
$F_k(\partial \mathcal{H}_a)$ in $M$ with the metric $g_{ij}(t)$
for $t$ close to $t_k$ and with the same area for each component
since $\partial \mathcal{H}_a$ is strictly concave and a
neighborhood of $\partial\mathcal{H}_a$ is foliated by constant
mean curvature hypersurfaces where the area $a$ has a nonzero
gradient and
$F_k:(\mathcal{H}_a,h_{ij})\rightarrow(M,g_{ij}(t_k))$ is close
enough to an isometry and $g_{ij}(t)$ varies smoothly. Then by
applying Lemma 5.3.8 we can smoothly continue the harmonic
diffeomorphism $F_k$ forward in time a little to a family of
harmonic diffeomorphisms $F_k(t)$ from $\mathcal{H}_a$ into $M$
with the metric $g_{ij}(t)$, with $F_k(t_k)=F_k$, where each
$F_k(t)$ takes $\partial \mathcal{H}_a$ into the constant mean
curvature hypersurface we just found in $(M,g_{ij}(t))$ and
satisfies the free boundary condition, and also satisfies
$$
||F^*_k(t)g_{ij}(t)-h_{ij}||_{C^q(\mathcal{H}_a)}<\delta.
$$
We claim that for all sufficiently large $k$, we can smoothly
extend the harmonic diffeomorphism $F_k$ to the family harmonic
diffeomorphisms $F_k(t)$ with
$||F^*_k(t)g_{ij}(t)-h_{ij}||_{C^q(\mathcal{H}_a)}\le\delta$ on a
maximal time interval $t_k\le t\leq\omega_k$ (or $t_k\leq
t<\omega_k$ when $\omega_k=+\infty$); and if $\omega_k<+\infty$,
then \begin{equation}
||F^*_k(\omega_k)g_{ij}(\omega_k)-h_{ij}||_{C^q(\mathcal{H}_a)}
=\delta.  %\eqno(5.3.9)
\end{equation}

Clearly the above argument shows that the set of $t$ where we can
extend the harmonic diffeomorphisms as desired is open. To verify
claim (5.3.9), we thus only need to show that if we have a family
of harmonic diffeomorphisms $F_k(t)$ such as we desire for $t_k\le
t<\omega (<+\infty)$, we can take the limit of $F_k(t)$ as
$t\rightarrow \omega$ to get a harmonic diffeomorphism
$F_k(\omega)$ satisfying
$$
||F^*_k(\omega)g_{ij}(\omega)-h_{ij}||_{C^q(\mathcal{H}_a)}\le\delta,
$$
and if
$$
||F^*_k(\omega)g_{ij}(\omega)-h_{ij}||_{C^q(\mathcal{H}_a)} <
\delta,
$$
then we can extend $F_k(\omega)$ forward in time a little (i.e.,
we can find a constant mean curvature hypersurface near
$F_k(\omega)(\partial \mathcal{H}_a)$ in $M$ with the metric
$g_{ij}(t)$ for each $t$ close to $\omega$ and with the same area
$a$ for each component). Note that \begin{equation}
||F^*_k(t)g_{ij}(t)-h_{ij}||_{C^q(\mathcal{H}_a)}<\delta.
%\eqno(5.3.10)
\end{equation} for $t_k\le t<\omega$ and the metrics $g_{ij}(t)$ for $t_k\le
t\le\omega$ are uniformly equivalent. We can find a subsequence
$t_n\rightarrow \omega$ for which $F_k(t_n)$ converge to
$F_k(\omega)$ in $C^{q-1}(\mathcal{H}_a)$ and the limit map has
$$
||F^*_k(\omega)g_{ij}(\omega)-h_{ij}||_{C^{q-1}(\mathcal{H}_a)}
\le\delta.
$$
We need to check that $F_k(\omega)$ is still a diffeomorphism. We
at least know $F_k(\omega)$ is a local diffeomorphism, and
$F_k(\omega)$ is the limit of diffeomorphisms, so the only
possibility of overlap is at the boundary. Hence we use the fact
that $F_k(\omega)(\partial\mathcal{H}_a)$ is still strictly
concave since $q$ is large and $\delta$ is small to prevent the
boundary from touching itself. Thus $F_k(\omega)$ is a
diffeomorphism. A limit of harmonic maps is harmonic, so
$F_k(\omega)$ is a harmonic diffeomorphism from $\mathcal{H}_a$
into $M$ with the metric $g_{ij}(\omega)$. Moreover $F_k(\omega)$
takes $\partial\mathcal{H}_a$ to the constant mean curvature
hypersurface $\partial (F_k(\omega)(\mathcal{H}_a))$ of the area
$a$ in $(M,g_{ij}(\omega))$ and continue to satisfy the free
boundary condition. As a consequence of the standard regularity
result of elliptic partial differential equations (see for example
\cite{GT}), the map $F_k(\omega)\in C^\infty(\mathcal{H}_a)$ and
then from (5.3.10) we have
$$
||F^*_k(\omega)g_{ij}(\omega)-h_{ij}||_{C^q(\mathcal{H}_a)}\le\delta.
$$
If
$||F^*_k(\omega)g_{ij}(\omega)-h_{ij}||_{C^q(\mathcal{H}_a)}=\delta$,
we then finish the proof of the claim. So we may assume that
$||F^*_k(\omega)g_{ij}(\omega)-h_{ij}||_{C^q(\mathcal{H}_a)}<\delta$.
We want to show that $F_k(\omega)$ can be extended forward in time
a little.

We argue by contradiction. Suppose not, then we consider the new
sequence of the manifolds $M$ with metric $g_{ij}(\omega)$ and the
origins $F_k(\omega)(P)$. Since $F_k(\omega)$ are close to
isometries, the injectivity radii of the metrics $g_{ij}(\omega)$
at $F_k(\omega)(P)$ do not go to zero, and we can extract a
subsequence which converges to a hyperbolic limit
$\widetilde{\mathcal{H}}$ with the metric $\widetilde{h}_{ij}$ and
the origin $\widetilde{P}$ and with finite volume. The new limit
$\widetilde{\mathcal{H}}$ has at least as many cusps as the old
limit $\mathcal{H}$, since we choose $\mathcal{H}$ with cusps as
few as possible. By the definition of convergence, we can find a
sequence of compact sets $\widetilde{B}_k$ exhausting
$\widetilde{\mathcal{H}}$ and containing $\widetilde{P}$, and a
sequence of diffeomorphisms $\widetilde{F}_k$ of neighborhoods of
$\widetilde{B}_k$ into $M$ with $\widetilde{F}_k(\widetilde{P}) =
F_k(\omega)(P)$ such that for each compact set $\widetilde{B}$ in
$\widetilde{\mathcal{H}}$ and each integer $m$
$$
||\widetilde{F}^*_k(g_{ij}(\omega)) -
\widetilde{h}_{ij}||_{C^m(\widetilde{B})} \rightarrow 0
$$
as $k \rightarrow +\infty$. For large enough $k$ the set
$\widetilde{F}_k(\widetilde{B}_k)$ will contain all points out to
any fixed distance we need from the point $F_k(\omega)(P)$, and
then
$$
\widetilde{F}_k(\widetilde{B}_k)\supset F_k(\omega)(\mathcal{H}_a)
$$
since the points of $\mathcal{H}_a$ have a bounded distance from
$P$ and $F_k(\omega)$ are reasonably close to preserving the
metrics. Hence we can form the composition
$$
G_k=\widetilde{F}^{-1}_k\circ F_k(\omega): \mathcal{H}_a
\rightarrow \widetilde{\mathcal{H}}.
$$
Arbitrarily fix $\delta' \in (\delta, \delta_0)$. Since the
$\widetilde{F}_k$ are as close to preserving the metric as we
like, we have
$$
||G^*_k\widetilde{h}_{ij} - h_{ij}||_{C^q(\mathcal{H}_a)} <
\delta'
$$
for all sufficiently large $k$. By Lemma 5.3.7, we deduce that
there exists an isometry $I$ of $\mathcal{H}$ to
$\widetilde{\mathcal{H}}$, and then
$(M,g_{ij}(\omega),F_k(\omega)(P))$ (on compact subsets) is very
close to $(\mathcal{H},h_{ij},P)$ as long as $\delta$ small enough
and $k$ large enough. Since $F_k(\omega)(\partial \mathcal{H}_a)$
is strictly concave and the foliation of a neighborhood of
$F_k(\omega)(\partial \mathcal{H}_a)$ by constant mean curvature
hypersurfaces has the area as a function with nonzero gradient, by
the implicit function theorem, there exists a unique constant mean
curvature hypersurface with the same area $a$ near
$F_k(\omega)(\partial \mathcal{H}_a)$ in $M$ with the metric
$g_{ij}(t)$ for $t$ close to $\omega$. Hence, when $k$
sufficiently large, $F_k(\omega)$ can be extended forward in time
a little. This is a contradiction and we have proved claim
(5.3.9).

We further claim that there must be some $k$ such that
$\omega_k=+\infty$ (i.e., we can smoothly continue the family of
harmonic diffeomorphisms $F_k(t)$ for all $t_k\le t<+\infty$, in
other words, there must be at least one hyperbolic piece
persisting). We argue by contradiction. Suppose for each $k$ large
enough, we can continue the family $F_k(t)$ for $t_k\le t\le
\omega_k<+\infty$ with
$$
||F^*_k(\omega_k)g_{ij}(\omega_k)-h_{ij}||_{C^q(\mathcal{H}_a)}
=\delta.
$$
Then as before, we consider the new sequence of the manifolds $M$
with metrics $g_{ij}(\omega_k)$ and origins $F_k(\omega_k)(P)$.
For sufficiently large $k$, we can obtain diffeomorphisms
$\widetilde{F}_k$ of neighborhoods of $\widetilde{B}_k$ into $M$
with $\widetilde{F}_k(\widetilde{P})=F_k(\omega_k)(P)$ which are
as close to preserving the metric as we like, where
$\widetilde{B}_k$ is a sequence of compact sets, exhausting some
hyperbolic three-manifold $\tilde{\mathcal{H}}$, of finite volume
and with no fewer cusps (than $\mathcal{H}$), and containing
$\tilde{P}$; moreover, the set $\widetilde{F}_k(\widetilde{B}_k)$
will contain all the points out to any fixed distance we need from
the point $F_k(\omega_k)(P);$ and hence
$$
\widetilde{F}_k(\widetilde{B}_k)\supseteq
F_k(\omega_k)(\mathcal{H}_a)
$$
since $\mathcal{H}_a$ is at bounded distance from $P$ and
$F_k(\omega_k)$ is reasonably close to preserving the metrics.
Then we can form the composition
$$
G_k=\widetilde{F}^{-1}_k\circ F_k(\omega_k):\quad
\mathcal{H}_a\rightarrow\tilde{\mathcal{H}}.
$$
Since the $\widetilde{F}_k$ are as close to preserving the metric
as we like, for any $\widetilde{\delta}>\delta$ we have
$$
||G^*_k\widetilde{h}_{ij}-h_{ij}||_{C^q(\mathcal{H}_a)}
<\widetilde{\delta}
$$
for large enough $k$. Then a subsequence of $G_k$ converges at
least in $C^{q-1}(\mathcal{H}_a)$ topology to a map $G_\infty$ of
$\mathcal{H}_a$ into $\widetilde{\mathcal{H}}$. By the same reason
as in the argument of previous two paragraphs, the limit map
$G_\infty$ is a smooth harmonic diffeomorphism from
$\mathcal{H}_a$ into $\widetilde{\mathcal{H}}$ with the metric
$\tilde{h}_{ij}$, and takes $\partial \mathcal{H}_a$ to a constant
mean curvature hypersurface $G_\infty(\partial \mathcal{H}_a)$ of
$(\tilde{\mathcal{H}},\widetilde{h}_{ij})$ with the area $a$, and
also satisfies the free boundary condition. Moreover we still have
\begin{equation} ||G^*_\infty\widetilde{h}_{ij}-h_{ij}||_{C^q(\mathcal{H}_a)}
=\delta. %\eqno(5.3.11)
\end{equation} Now by Lemma 5.3.7 we deduce that there exists an isometry $I$
of $\mathcal{H}$ to $\widetilde{\mathcal{H}}$ with
$$
d_{C^{l_0}(\mathcal{H}_a)}(G_\infty,I)<\varepsilon_0.
$$
By using $I$ to identify $\widetilde{\mathcal{H}}_a$ and
$\mathcal{H}_a$, we see that the map $I^{-1}\circ G_\infty $ is a
harmonic diffeomorphism of $\mathcal{H}_a$ to itself which
satisfies the free boundary condition and
$$
d_{C^{l_0}(\mathcal{H}_a)}(I^{-1}\circ G_\infty,Id)<\varepsilon_0.
$$
{}From the uniqueness in Lemma 5.3.8 we conclude that $I^{-1}\circ
G_\infty =Id$ which contradicts with (5.3.11). This shows at least
one hyperbolic piece persists. Moreover the pull-back of the
solution metric $g_{ij}(t)$ by $F_k(t)$, for $t_k\leq t<+\infty$,
is as close to the hyperbolic metric $h_{ij}$ as we like.

We can continue to form other persistent hyperbolic pieces in the
same way as long as there are any points $P_k$ outside of the
chosen pieces where the injectivity radius at times
$t_k\rightarrow \infty$ are all at least some fixed positive
number $\rho >0$. The only modification in the proof is to take
the new limit $\mathcal{H}$ to have the least possible number of
cusps out of all remaining possible limits.

Note that the volume of the normalized Ricci flow is constant in
time. Therefore by combining with Margulis lemma (see for example
\cite{Gro79} \cite{KM}), we have proved that there exists a finite
collection of complete noncompact hyperbolic three-manifolds
${\mathcal{H}}_1,\ldots,{\mathcal{H}}_m $ with finite volume, a
small number $a>0$ and a time $T<+\infty$ such that for all $t$
beyond $T$ we can find diffeomorphisms $\varphi_l(t)$ of
$({\mathcal{H}}_{l})_a $ into $M$, $1\leq l\leq m$, so that the
pull-back of the solution metric $g_{ij}(t)$ by $\varphi_l(t)$ is
as close to the hyperbolic metrics as we like and the exceptional
part of $M$ where the points are not in the image of any
$\varphi_l$ has the injectivity radii everywhere as small as we
like.

\medskip
{\bf Part III: Incompressibility} \vskip 0.2cm

We remain to show that the boundary tori of any persistent
hyperbolic piece are incompressible, in the sense that the
fundamental group of the torus injects into that of the whole
manifold. The argument of this part is a parabolic version of
Schoen and Yau's minimal surface argument in \cite{ScY79m, ScY79,
ScY82}.

Let $B$ be a small positive number and assume the above positive
number $a$ is much smaller than $B$. Denote by $M_a$ a persistent
hyperbolic piece of the manifold $M$ truncated by boundary tori of
area $a$ with constant mean curvature and denote by $M_a^c=M
\setminus \stackrel{\circ}{M_a}$ the part of $M$ exterior to
$M_a$. Thus there is a persistent hyperbolic piece $M_B \subset
M_a$ of the manifold $M$ truncated by boundary tori of area $B$
with constant mean curvature. We also denote by $M_B^c=M \setminus
\stackrel{\circ}{M_B}$. By Van Kampen's Theorem, if
$\pi_1(\partial M_B)$ injects into $\pi_1(M_B^c)$ then it injects
into $\pi_1(M)$ also. Thus we only need to show $\pi_1(\partial
M_B)$ injects into $\pi_1(M_B^c)$.

We will argue by contradiction. Let $T$ be a torus in $\partial
M_B$. Suppose $\pi_1(T)$ does not inject into $\pi_1(M_B^c)$, then
by Dehn's Lemma the kernel is a cyclic subgroup of $\pi_1(T)$
generated by a primitive element. The work of Meeks-Yau \cite{MY}
or Meeks-Simon-Yau \cite{MSY} shows that among all disks in
$M_B^c$ whose boundary curve lies in $T$ and generates the kernel,
there is a smooth embedded disk normal to the boundary which has
the least possible area. Let $A=A(t)$ be the area of this disk.
This is defined for all $t$ sufficiently large. We will show that
$A(t)$ decreases at a rate bounded away from zero which will be a
contradiction.

Let us compute the rate at which $A(t)$ changes under the Ricci
flow. We need to show $A(t)$ decrease at least at a certain rate,
and since $A(t)$ is the minimum area to bound any disk in the
given homotopy class, it suffices to find some such disk whose
area decreases at least that fast. We choose this disk as follows.
Pick the minimal disk  at time $t_0$, and extend it smoothly a
little past the boundary torus since the minimal disk is normal to
the boundary. For times $t$ a little bigger than $t_0$, the
boundary torus may need to move a little to stay constant mean
curvature with area $B$ as the metrics change, but we leave the
surface alone and take the bounding disk to be the one cut off
from it by the new torus. The change of the area $\tilde{A}(t)$ of
such disk comes from the change in the metric and the change in
the boundary.

For the change in the metric, we choose an orthonormal frame
$X,Y,Z$ at a point $x$ in the disk so that $X$ and $Y$ are tangent
to the disk while $Z$ is normal and compute the rate of change of
the area element $d\sigma$ on the disk as
\begin{align*}
\frac{\partial}{\partial t}d\sigma&  =
\frac{1}{2}(g^{ij})^T\left(\frac{2}{3}r(g_{ij}\right)^T-2(R_{ij})^T)d\sigma\\[2mm]
&  = \left[\frac{2}{3}r-\Ric(X,X)-\Ric(Y,Y)\right]d\sigma,
\end{align*}
since the metric evolves by the normalized Ricci flow. Here
$(\cdot)^T$ denotes the tangential projections on the disk. Notice
the torus $T$ may move in time to preserve constant mean curvature
and constant area $B$. Suppose the boundary of the disk evolves
with a normal velocity $N$. The change of the area at boundary
along a piece of length $ds$ is given by $Nds$. Thus the total
change of the area $\tilde{A}(t)$ is given by
$$
\frac{d\tilde{A}}{dt}
=\int\int\left(\frac{2}{3}r-\Ric(X,X)-\Ric(Y,Y)\right)d\sigma
+\int_\partial Nds.
$$
Note that
\begin{align*}
\Ric(X,X)+\Ric(Y,Y)&  = R(X,Y,X,Y)+R(X,Z,X,Z)\\
&\quad +R(Y,X,Y,X)+R(Y,Z,Y,Z)\\
&  = \frac{1}{2}R+R(X,Y,X,Y).
\end{align*}
By the Gauss equation, the Gauss curvature $K$ of the disk is
given by
$$
K=R(X,Y,X,Y)+\det II
$$
where $II$ is the second fundamental form of the disk in $M_B^c$.
This gives at $t=t_0$,
$$
\frac{dA}{dt}\leq\int\int\left(\frac{2}{3}r-\frac{1}{2}R\right)d\sigma
-\int\int(K-\det II)d\sigma+\int_\partial N ds
$$
Since the bounding disk is a minimal surface, we have
$$
\det II\leq 0.
$$
The Gauss-Bonnet Theorem tells us that for a disk
$$
\int\int Kd\sigma+\int_\partial k ds=2\pi
$$
where $k$ is the geodesic curvature of the boundary. Thus we
obtain \begin{equation} \frac{dA}{dt}\leq
\int\int\left(\frac{2}{3}r-\frac{1}{2}R\right)d\sigma+\int_\partial
kds+\int_\partial Nds -2\pi. %\eqno(5.3.12)
\end{equation} Recall that we are assuming $R_{min}(t)$ increases
monotonically to $-6$ as $t\rightarrow +\infty$. By the evolution
equation of the scalar curvature,
$$
\frac{d}{dt}R_{\min}(t)\geq 4(r(t)-R_{\min}(t))
$$
and then
$$
\int_0^\infty (r(t)-R_{\min}(t))dt<+\infty.
$$
This implies that $r(t)\rightarrow -6$ as $t\rightarrow +\infty$
by using the derivatives estimate for the curvatures. Thus for
every $\varepsilon >0$ we have
$$
\frac{2}{3}r-\frac{1}{2}R\leq -(1-\varepsilon)
$$
for $t$ sufficiently large. And then the first term on RHS of
(5.3.12) is bounded above by
$$
\int\int \left(\frac{2}{3}r-\frac{1}{2}R\right)d\sigma\leq
-(1-\varepsilon)A.
$$
The geodesic curvature $k$ of the boundary of the minimal disk is
the acceleration of a curve moving with unit speed along the
intersection of the disk with the torus; since the disk and torus
are normal, this is the same as the second fundamental form of the
torus in the direction of the curve of intersection. Now if the
metric were actually hyperbolic, the second fundamental form of
the torus would be exactly 1 in all directions. Note that the
persistent hyperbolic pieces are as close to the standard
hyperbolic as we like. This makes that the second term of RHS of
(5.3.12) is bounded above by
$$
\int_\partial kds\leq(1+\varepsilon_0)L
$$
for some sufficiently small positive number $\varepsilon_0>0$,
where $L$ is the length of the boundary curve. Also since the
metric on the persistent hyperbolic pieces are close to the
standard hyperbolic as we like, its change under the normalized
Ricci flow is as small as we like; So the motion of the constant
mean curvature torus of fixed area $B$ will have a normal velocity
$N$ as small as we like. This again makes the third term of RHS of
(5.3.12) bounded above by
$$
\int_\partial Nds\leq \varepsilon_0 L.
$$
Combining these estimates, we obtain \begin{equation}
\frac{dA}{dt}\leq(1+2\varepsilon_0)L-(1-\varepsilon_0)A-2\pi
%\eqno(5.3.13)
\end{equation} on the persistent hyperbolic piece, where $\varepsilon_0$ is
some sufficiently small positive number.

We next need to bound the length $L$ in terms of the area $A$.
Since $a$ is much smaller than $B$, for large $t$ the metric is as
close as we like to the standard hyperbolic one; not just on the
persistent hyperbolic piece $M_B$ but as far beyond as we like.
Thus for a long distance into $M_B^c$ the metric will look nearly
like a standard hyperbolic cusplike collar.

Let us first recall a special coordinate system on the standard
hyperbolic cusp projecting beyond torus $T_1$ in $\partial
{\mathcal {H}}_1$ as follows. The universal cover of the flat
torus $T_1$ can be mapped conformally to the $x$-$y$ plane so that
the deck transformation of $T_1$ become translations in $x$ and
$y$, and so that the Euclidean area of the quotient is 1; then
these coordinates are unique up to a translation. The hyperbolic
cusp projecting beyond the torus $T_1$ in $\partial {\mathcal
{H}}_1$  can be parametrized by $\{(x,y,z)\in {\mathbb{R}}^3\ |\ z
> 0\}$ with the hyperbolic metric
\begin{equation}
ds^2=\frac{dx^2+dy^2+dz^2}{z^2}. %\eqno(5.3.14)
\end{equation} Note that we can make the solution metric, in an arbitrarily
large neighborhood of the torus $T$ (of $\partial M_B$), as close
to hyperbolic as we wish (in the sense that there exists a
diffeomorphism from a large neighborhood of the torus $T_B$ (of
$\partial\mathcal{H}_B$) on the standard hyperbolic cusp to the
above neighborhood of the torus $T$ (of $\partial M_B$) such that
the pull-back of the solution metric by the diffeomorphism is as
close to the hyperbolic metric as we wish). Then by using this
diffeomorphism (up to a slight modification) we can parametrize
the cusplike tube of $M_B^c$ projecting beyond the torus $T$ in
$\partial M_B$ by $\{(x,y,z)\ |\  z\geq \zeta\}$ where the height
$\zeta$ is chosen so that the torus in the hyperbolic cusp at
height $\zeta$ has the area $B$.

Now consider our minimal disk, and let $L(z)$ be the length of the
curve of the intersection of the disk with the torus at height $z$
in the above coordinate system, and also let $A(z)$ be the area of
the part of the disk between $\zeta$ and $z$. We now want to
derive a monotonicity formula on the area $A(z)$ for the minimal
surface.

For almost every $z$ the intersection of the disk with the torus
at height $z$ is a smooth embedded curve or a finite union of them
by the standard transversality theorem. If there is more than one
curve, at least one of them is not homotopic to a point in $T$ and
represents the primitive generator in the kernel of $\pi_1(T)$
such that a part of the original disk beyond height $z$ continues
to a disk that bounds it. We extend this disk back to the initial
height $\zeta$ by dropping the curve straight down. Let
$\tilde{L}(z)$ be the length of the curve we picked at height $z$;
of course $\tilde{L}(z)\leq L(z)$ with equality if it is the only
piece. Let $\tilde{L}(w)$ denote the length of the same curve in
the $x$-$y$ plane dropped down to height $w$ for $\zeta\leq w \leq
z$. In the hyperbolic space we would have
$$
\tilde{L}(w)=\frac{z}{w}\tilde{L}(z)
$$
exactly. In our case there is a small error proportional to
$\tilde{L}(z)$ and we can also take it proportional to the
distance $z-w$ by which it drops since
$\tilde{L}(w)|_{w=z}=\tilde{L}(z)$ and the solution metric is
close to the hyperbolic in the $C_{loc}^{\infty}$ topology. Thus,
for arbitrarily given $\delta>0$ and $\zeta^{\ast}>\zeta$, as the
solution metric is sufficiently close to hyperbolic, we have
$$
|\tilde{L}(w)-\frac{z}{w}\tilde{L}(z)|\leq \delta(z-w)\tilde{L}(z)
$$
for all $z$ and $w$ in $\zeta\leq w\leq z\leq {\zeta}^{\ast}$. Now
given $\varepsilon$ and ${\zeta}^{\ast}$ pick
$\delta=2\varepsilon/\zeta^{\ast} $. Then \begin{equation} \tilde{L}(w)\leq
\frac{z}{w}\tilde{L}(z)\left[1+\frac{2\varepsilon(z-w)}{w}\right].
%\eqno(5.3.15)
\end{equation} When we drop the curve vertically for the construction of the
new disk we get an area $\tilde{A}(z)$ between $\zeta$ and $z$
given by
$$
\tilde{A}(z)=(1+o(1))\int_{\zeta}^{z} \frac{\tilde{L}(w)}{w}dw.
$$
Here and in the following $o(1)$ denotes various small error
quantities as the solution metric close to hyperbolic. On the
other hand if we do not drop vertically we pick up even more area,
so the area $A(z)$ of the original disk between $\zeta$ and $z$
has \begin{equation} A(z)\geq (1-o(1))\int_{\zeta}^{z} \frac{L(w)}{w} dw.
%\eqno(5.3.16)
\end{equation} Since the original disk minimized among all disks bounded a
curve in the primitive generator of the kernel of $\pi_1(T)$, and
the new disk beyond the height $z$ is part of the original disk,
we have
$$
A(z)\leq \tilde{A}(z)
$$
and then by combining with (5.3.15),
\begin{align*}
\int_{\zeta}^{z} \frac{L(w)}{w}dw &  \leq (1+o(1))
z\tilde{L}(z)\int_{\zeta}^{z}
\left[\frac{1-2\varepsilon}{w^2}+\frac{2\varepsilon z}{w^3}\right]dw \\
& \leq (1+o(1))L(z)\frac{(z-\zeta)}{\zeta}
\left[1+\varepsilon\left(\frac{z-\zeta}{\zeta}\right)\right].
\end{align*}
Here we used the fact that $\tilde{L}(z)\leq L(z)$. Since the
solution metric is sufficiently close to hyperbolic, we have
\begin{align*}
\frac{d}{dz}\left(\int_{\zeta}^{z}\frac{L(w)}{w}dw\right)
&  = \frac{L(z)}{z}\\
& \geq(1-o(1))\frac{\zeta}{z(z-\zeta)}
\left[1-\varepsilon\left(\frac{z-\zeta}{\zeta}\right)\right]
\int_{\zeta}^{z}\frac{L(w)}{w}dw \\
&  \geq\left[\frac{1}{z-\zeta}-\frac{1+2\varepsilon}{z}\right]
\int_{\zeta}^{z}\frac{L(w)}{w}dw,
\end{align*}
or equivalently \begin{equation} \frac{d}{dz}\log\left\{
\frac{z^{1+2\varepsilon}}{(z-\zeta)}
\int_{\zeta}^{z}\frac{L(w)}{w}dw\right\}\geq 0.
%\eqno(5.3.17)
\end{equation} This is the desired monotonicity formula for the area $A(z)$.

It follows directly from (5.3.16) and (5.3.17) that
$$
\frac{z^{1+2\varepsilon}}{(z-\zeta)}A(z)\geq
(1-o(1))\zeta^{2\varepsilon}L(\zeta),
$$
or equivalently
$$
L(\zeta)\leq (1+o(1))
\left(\frac{z}{\zeta}\right)^{2\varepsilon}\frac{z}{z-\zeta} A(z)
$$
for all $z\in [\zeta,\zeta^{\ast}]$. Since the solution metric, in
an arbitrarily large neighborhood of the torus $T$ (of $\partial
M_B)$, as close to hyperbolic as we wish, we may assume that
${\zeta}^{\ast}$ is so large that $\sqrt{\zeta^{\ast}}>\zeta$ and
$\frac{\sqrt{\zeta^{\ast}}}{\sqrt{\zeta^{\ast}}-\zeta}$ is close
to 1, and also $\varepsilon >0$ is so small that
$(\frac{\sqrt{\zeta^{\ast}}}{\zeta})^{2\varepsilon}$ is close to
1. Thus for arbitrarily small $\delta_0 > 0$, we have \begin{equation}
L(\zeta)\leq (1+\delta_0)A\left(\sqrt{\zeta^{\ast}}\right). %\eqno (5.3.18)
\end{equation} Now recall that (5.3.13) states
$$
\frac{dA}{dt}\leq (1+2\varepsilon_0)L-(1-\varepsilon_0)A-2\pi.
$$
We now claim that if
$$
(1+2\varepsilon_0)L-(1-\varepsilon_0)A\geq 0
$$
then $L=L(\zeta)$ is uniformly bounded from above.

Indeed by the assumption we have
$$
A(\zeta^{\ast})\leq
\frac{(1+2\varepsilon_0)}{(1-\varepsilon_0)}L(\zeta)
$$
since $A(\zeta^{\ast})\leq A$. By combining with (5.3.16) we have
some $z_0\in (\sqrt{\zeta^{\ast}},\zeta^{\ast})$ satisfying
\begin{align*}
\frac{L(z_0)}{z_0}\left(\zeta^{\ast}-\sqrt{\zeta^{\ast}}\right)&  \leq
\int_{\sqrt{\zeta^{\ast}}}^{\zeta^{\ast}} \frac{L(w)}{w}dw \\
&  \leq  (1+o(1))A(\zeta^{\ast}) \\
&  \leq
(1+o(1))\left(\frac{1+2\varepsilon_0}{1-\varepsilon_0}\right)L(\zeta).
\end{align*}
Thus for $\zeta^{\ast}$ suitably large, by noting that the
solution metric on a large neighborhood of $T$ (of $\partial
{M}_B)$ is sufficiently close to hyperbolic, we have \begin{equation}
L(z_0)\leq (1+4\varepsilon_0)\frac{z_0}{\zeta^{\ast}}L(\zeta)
%\eqno (5.3.19)
\end{equation} for some $z_0\in (\sqrt{\zeta^{\ast}},\zeta^{\ast})$. It is
clear that we may assume the intersection curve between the
minimal disk with the torus at this height $z_0$ is smooth and
embedded. If the intersection curve at the height $z_0$ has more
than one piece, as before one of them will represent the primitive
generator in the kernel of $\pi_1(T)$, and we can ignore the
others. Let us move (the piece of) the intersection curve on the
torus at height $z_0$ through as small as possible area in the
same homotopy class of $\pi_1(T)$ to a curve which is a geodesic
circle in the flat torus coming from our special coordinates, and
then drop this geodesic circle vertically in the special
coordinates to obtain another new disk. We will compare the area
of this new disk with the original minimal disk as follows.

Denote by $G$ the length of the geodesic circle in the standard
hyperbolic cusp at height 1. Then the length of the geodesic
circle at height $z_0$ will be $G/z_0$. Observe that given an
embedded curve of length $l$ circling the cylinder $S^1\times
\mathbb{R}$ of circumference $w$ once, it is possible to deform
the curve through an area not bigger than $lw$ into a meridian
circle. Note that (the piece of) the intersection curve represents
the primitive generator in the kernel of $\pi_1(T)$. Note also
that the solution metric is sufficiently close to the hyperbolic
metric. Then the area of the deformation from (the piece of) the
intersection curve on the torus at height $z_0$ to the geodesic
circle at height $z_0$ is bounded by
$$
(1+o(1))\left(\frac{G}{z_0}\right)\cdot L(z_0).
$$
The area to drop the geodesic circle from height $z_0$ to height
$\zeta$ is bounded by
$$
(1+o(1))\int_{\zeta}^{z_0}\frac{G}{w^2}dw.
$$
Hence comparing the area of the original minimal disk to that of
this new disk gives
$$
A(z_0)\leq (1+o(1))G\left[\frac{L(z_0)}{z_0}
+\left(\frac{1}{\zeta}-\frac{1}{z_0}\right)\right].
$$
By (5.3.18), (5.3.19) and the fact that
$z_0\in(\sqrt{\zeta^{\ast}},\zeta^{\ast})$, this in turn gives
\begin{align*}
L(\zeta)&  \leq (1+\delta_0)A(z_0) \\
&  \leq(1+\delta_0)G
\left[(1+4\varepsilon_0)\frac{L(\zeta)}{\zeta^{\ast}}+\frac{1}{\zeta}\right].
\end{align*}
Since $\zeta^{\ast}$ is suitably large, we obtain
$$
L(\zeta)\leq 2G/\zeta
$$
This gives the desired assertion since $G$ is fixed from the
geometry of the limit hyperbolic manifold $\mathcal{H}$ and
$\zeta$ is very large as long as the area $B$ of $\partial M_B$
small enough.

Thus the combination of (5.3.13), (5.3.18) and the assertion
implies that either
$$
\frac{d}{dt}A\leq-2\pi,
$$
or
\begin{align*}
\frac{d}{dt}A&  \leq (1+2\varepsilon_0)L-(1-\varepsilon_0)A-2\pi\\
&  \leq (1+2\varepsilon_0)\frac{2G}{\zeta} - 2\pi\\
&  \leq  -\pi,
\end{align*}
since the solution metric on a very large neighborhood of the
torus $T$ (of $\partial {M}_B$) is sufficiently close to
hyperbolic and $\zeta$ is very large as the area $B$ of $\partial
M_B$ small enough. This is impossible because $A\geq 0$ and the
persistent hyperbolic pieces go on forever. The contradiction
shows that $\pi_1(T)$ in fact injects into $\pi_1(M_B^c)$. This
proves that $\pi_1(\partial{M}_B) $ injects into $\pi_1(M)$.

Therefore we have completed the proof of Theorem 5.3.4.
%$$\eqno\#$$
\endproof


%\bigskip
\newpage

\section{Named intermediate formalization obligations}

\begin{lemma}[curvature-gradient estimate under three-dimensional pinching]
  \label{lem:cz-three-dimensional-pinching-curvature-gradient}
  \dcref{Claim in the proof of Proposition 5.2.5}
  \group{Differentiable Sphere Theorems}
  \source{cao-zhu}
For any $\varepsilon>0$, there exists a positive
constant $C_\varepsilon<+\infty$ such that for any time $\tau\ge
0$, we have
$$
\max_{t\le\tau}\max_{x\in M}|\nabla Rm(x,t)| \le \varepsilon
\max_{t\le\tau}\max_{x\in M}|Rm(x,t)|^\frac{3}{2}+C_\varepsilon.
$$

\end{lemma}
\begin{proof}
We argue by contradiction. Suppose the above gradient estimate
fails for some fixed $\varepsilon_0>0$. Pick a sequence
$C_j\rightarrow+\infty$, and pick points $x_j\in M$ and times
$\tau_j$ such that
$$
|\nabla Rm(x_j,\tau_j)|\ge\varepsilon_0\max_{t\le\tau_j}
\max_{x\in M}|Rm(x,t)|^\frac{3}{2}+C_j,\quad j=1,2,\ldots.
$$
Choose $x_j$ to be the origin, and pull the metric back to a small
ball on the tangent space $T_{x_j}M$ of radius $r_j$ proportional
to the reciprocal of the square root of the maximum curvature up
to time $\tau_j$ (i.e., $\max_{t\le\tau_j}\max_{x\in M}|Rm(x,$
$t)|$). Clearly the maximum curvatures go to infinity by Shi's
derivative estimate of curvature (Theorem 1.4.1). Dilate the
metrics so that the maximum curvature
$$
\max_{t\le\tau_j}\max_{x\in M}|Rm(x,t)|
$$
becomes 1 and translate time so that $\tau_j$ becomes the time 0.
By Theorem 4.1.5, we can take a (local) limit. The limit metric
satisfies
$$
|\nabla Rm(0,0)|\ge \varepsilon_0>0.
$$
However the pinching estimate in Proposition 5.2.5 tells us the
limit metric has
$$
R_{ij}-\frac{1}{3}Rg_{ij}\equiv 0.
$$
By using the contracted second Bianchi identity
$$
\frac{1}{2}\nabla_iR =\nabla^jR_{ij}
=\nabla^j\left(R_{ij}-\frac{1}{3}Rg_{ij}\right)+\frac{1}{3}\nabla_iR,
$$
we get
$$
\nabla_iR\equiv0\quad \text{and then}\quad \nabla_iR_{jk}\equiv0.
$$
For a three-manifold, this in turn implies
$$
\nabla Rm=0
$$
which is a contradiction. Hence we have proved the gradient
estimate claimed.

\end{proof}

\begin{lemma}[bounded cusp length under nonnegative area derivative]
  \label{lem:cz-hyperbolic-cusp-length-bounded}
  \dcref{Claim in the proof of Lemma 5.3.8}
  \group{Nonsingular Three-Dimensional Solutions}
  \source{cao-zhu}
If
$$
(1+2\varepsilon_0)L-(1-\varepsilon_0)A\geq 0
$$
then $L=L(\zeta)$ is uniformly bounded from above.

\end{lemma}
\begin{proof}
$$
A(\zeta^{\ast})\leq
\frac{(1+2\varepsilon_0)}{(1-\varepsilon_0)}L(\zeta)
$$
since $A(\zeta^{\ast})\leq A$. By combining with (5.3.16) we have
some $z_0\in (\sqrt{\zeta^{\ast}},\zeta^{\ast})$ satisfying
\begin{align*}
\frac{L(z_0)}{z_0}\left(\zeta^{\ast}-\sqrt{\zeta^{\ast}}\right)&  \leq
\int_{\sqrt{\zeta^{\ast}}}^{\zeta^{\ast}} \frac{L(w)}{w}dw \\
&  \leq  (1+o(1))A(\zeta^{\ast}) \\
&  \leq
(1+o(1))\left(\frac{1+2\varepsilon_0}{1-\varepsilon_0}\right)L(\zeta).
\end{align*}
Thus for $\zeta^{\ast}$ suitably large, by noting that the
solution metric on a large neighborhood of $T$ (of $\partial
{M}_B)$ is sufficiently close to hyperbolic, we have \begin{equation}
L(z_0)\leq (1+4\varepsilon_0)\frac{z_0}{\zeta^{\ast}}L(\zeta)
%\eqno (5.3.19)
\end{equation} for some $z_0\in (\sqrt{\zeta^{\ast}},\zeta^{\ast})$. It is
clear that we may assume the intersection curve between the
minimal disk with the torus at this height $z_0$ is smooth and
embedded. If the intersection curve at the height $z_0$ has more
than one piece, as before one of them will represent the primitive
generator in the kernel of $\pi_1(T)$, and we can ignore the
others. Let us move (the piece of) the intersection curve on the
torus at height $z_0$ through as small as possible area in the
same homotopy class of $\pi_1(T)$ to a curve which is a geodesic
circle in the flat torus coming from our special coordinates, and
then drop this geodesic circle vertically in the special
coordinates to obtain another new disk. We will compare the area
of this new disk with the original minimal disk as follows.

Denote by $G$ the length of the geodesic circle in the standard
hyperbolic cusp at height 1. Then the length of the geodesic
circle at height $z_0$ will be $G/z_0$. Observe that given an
embedded curve of length $l$ circling the cylinder $S^1\times
\mathbb{R}$ of circumference $w$ once, it is possible to deform
the curve through an area not bigger than $lw$ into a meridian
circle. Note that (the piece of) the intersection curve represents
the primitive generator in the kernel of $\pi_1(T)$. Note also
that the solution metric is sufficiently close to the hyperbolic
metric. Then the area of the deformation from (the piece of) the
intersection curve on the torus at height $z_0$ to the geodesic
circle at height $z_0$ is bounded by
$$
(1+o(1))\left(\frac{G}{z_0}\right)\cdot L(z_0).
$$
The area to drop the geodesic circle from height $z_0$ to height
$\zeta$ is bounded by
$$
(1+o(1))\int_{\zeta}^{z_0}\frac{G}{w^2}dw.
$$
Hence comparing the area of the original minimal disk to that of
this new disk gives
$$
A(z_0)\leq (1+o(1))G\left[\frac{L(z_0)}{z_0}
+\left(\frac{1}{\zeta}-\frac{1}{z_0}\right)\right].
$$
By (5.3.18), (5.3.19) and the fact that
$z_0\in(\sqrt{\zeta^{\ast}},\zeta^{\ast})$, this in turn gives
\begin{align*}
L(\zeta)&  \leq (1+\delta_0)A(z_0) \\
&  \leq(1+\delta_0)G
\left[(1+4\varepsilon_0)\frac{L(\zeta)}{\zeta^{\ast}}+\frac{1}{\zeta}\right].
\end{align*}
Since $\zeta^{\ast}$ is suitably large, we obtain
$$
L(\zeta)\leq 2G/\zeta
$$
This gives the desired assertion since $G$ is fixed from the
geometry of the limit hyperbolic manifold $\mathcal{H}$ and
$\zeta$ is very large as long as the area $B$ of $\partial M_B$
small enough.

Thus the combination of (5.3.13), (5.3.18) and the assertion
\end{proof}
